\chapter{Finite ramification and Herbrand's theorem}

Fix a prime $p$. A local field here is a complete discretely valued field
with perfect residue field of characteristic $p$; its residue field need
not be finite. Write $\mathcal O_E$, $\mathfrak m_E$, and $k_E$ for the
valuation ring, maximal ideal, and residue field of $E$. Normalize
$v_E(E^\times)=\mathbb Z$ and set $v_E(0)=+\infty$.
For a finite extension $E/K$, the extended valuation satisfies
$v_E|_K=e(E/K)v_K$. We use the basic structure theorem for finite
extensions of complete discretely valued fields: the valuation extends
uniquely, $\mathcal O_E$ is a finite free $\mathcal O_K$-module,
$[E:K]=e(E/K)[k_E:k_K]$, and ramification indices multiply in towers.
We also use Hensel's lemma for a simple residue root.
These are the local algebra foundations of \cite[Chapters I--II]{Serre1979}.
All extensions below are separable and lie in a fixed separable closure
$K^{\mathrm{sep}}$. Put $G_K=\Gal(K^{\mathrm{sep}}/K)$ and
$G_L=\Gal(K^{\mathrm{sep}}/L)$.

The finite arguments follow \cite[Chapter IV, Sections 1 and 3]{Serre1979}.
We write $j_E$ for Serre's unshifted order function; the number $i_E$ used
by \cite{Wintenberger1983} is $j_E-1$.

\section{Integral generators and lower groups}

\begin{lemma}[Integer congruence groups]
\label{fw:integer-lower-groups}
\dcref{Serre1979, Chapter IV, Section 1, Proposition 1}
\source{serre-1979-local-fields:page-0068,serre-1979-local-fields:page-0069}
\lean{FontaineWintenberger.lowerRamificationGroup, FontaineWintenberger.mem_lowerRamificationGroup_iff, FontaineWintenberger.lowerRamificationGroup_neg_one, FontaineWintenberger.lowerRamificationGroup_of_le_neg_one, FontaineWintenberger.lowerRamificationGroup_zero, FontaineWintenberger.lowerRamificationGroup_nat, FontaineWintenberger.lowerRamificationGroup_antitone, FontaineWintenberger.lowerRamificationGroup_subgroupOf, FontaineWintenberger.lowerRamificationGroup_map_subtype, FontaineWintenberger.lowerRamificationGroup_normal, FontaineWintenberger.smul_maximalIdeal, FontaineWintenberger.maximalIdeal_lowerRamificationGroup_normal, FontaineWintenberger.lowerRamificationGroup_zero_eq_ker_residue, ValuationSubring.decompositionSubgroupSMul, ValuationSubring.decompositionSubgroup_faithfulSMul, ValuationSubring.lowerRamificationGroup, ValuationSubring.mem_lowerRamificationGroup_iff, ValuationSubring.lowerRamificationGroup_neg_one, ValuationSubring.lowerRamificationGroup_zero, ValuationSubring.lowerRamificationGroup_antitone, ValuationSubring.lowerRamificationGroup_normal}
\leanok
Let a group $D$ act by ring automorphisms on a commutative local ring $A$,
with maximal ideal $\mathfrak m$. For $n\in\mathbb Z$ put
\[
 D_n=\{s\in D:\ s(a)-a\in\mathfrak m^{\max(n+1,0)}
                      \text{ for every }a\in A\}.
\]
These are decreasing normal subgroups, $D_n=D$ for $n\leq-1$, and
$D_0$ is the kernel of the action on the residue field.
For every subgroup $H\subset D$, the same construction for its restricted
action gives $H_n=H\cap D_n$. In particular, this construction applies
to a valuation subring of a field and its decomposition group, whose
action on the valuation subring is faithful.
This is the ideal-theoretic part of Serre's argument, extended to arbitrary
commutative local rings; eventual triviality requires the additional
discreteness and finiteness hypotheses considered below.
\end{lemma}
\begin{proof}
For any ideal $I$, the set of automorphisms acting trivially modulo $I$
is a subgroup: the identity has zero difference,
\[
 st(a)-a=(s(t(a))-t(a))+(t(a)-a),\qquad
 s^{-1}(a)-a=-(s(s^{-1}(a))-s^{-1}(a)).
\]
These identities prove closure under multiplication and inverse.
Every ring automorphism of a local ring preserves its maximal ideal,
since it preserves units, and therefore preserves each power of that
ideal. If $s$ acts trivially modulo such a power $I$, then
\[
 rsr^{-1}(a)-a=r\bigl(s(r^{-1}(a))-r^{-1}(a)\bigr)\in I,
\]
which proves normality. Powers of an ideal decrease with the exponent,
which gives the asserted monotonicity. At $n\leq-1$ the ideal is $A$.
At $n=0$ the difference condition is precisely triviality on $A/\mathfrak m$;
every residue class has a representative in $A$, so it is also sufficient.
Restricting the acting group leaves the condition on each element
unchanged, proving the subgroup identity.
For the last faithfulness assertion, every field element or its inverse
lies in the valuation subring. Thus two automorphisms agreeing on that
subring agree on every element of the field.
\end{proof}

\begin{lemma}[Ideal powers and the integral valuation]
\label{fw:ideal-power-valuation}
\dcref{Serre1979, Chapter I, Section 1; Chapter IV, Section 1}
\source{serre-1979-local-fields:page-0068}
\uses{fw:integer-lower-groups}
\lean{FontaineWintenberger.mem_maximalIdeal_pow_iff_addVal, FontaineWintenberger.addVal_smul, FontaineWintenberger.mem_lowerRamificationGroup_iff_addVal, FontaineWintenberger.mem_lowerRamificationGroup_nat_iff_addVal, ValuationSubring.mem_lowerRamificationGroup_iff_addVal}
\leanok
Let $A$ be a discrete valuation ring with maximal ideal $\mathfrak m$
and normalized additive valuation $v_A:A\to\mathbb N\cup\{+\infty\}$.
For $a\in A$ and $n\geq0$,
\[
 a\in\mathfrak m^n\quad\Longleftrightarrow\quad n\leq v_A(a).
\]
Every ring automorphism of $A$ preserves $v_A$.
Consequently the integer congruence condition defining $D_n$ is
$v_A(s(a)-a)\geq\max(n+1,0)$ for every $a\in A$.
\end{lemma}
\begin{proof}
Choose a uniformizer $\pi$. Then $\mathfrak m^n=(\pi^n)$,
so the membership assertion is equivalent to $\pi^n\mid a$.
The factorization of each nonzero element as a unit times a power of
$\pi$ gives $b\mid a$ if and only if $v_A(b)\leq v_A(a)$;
this also holds when either element is zero, with $v_A(0)=+\infty$.
Since $v_A(\pi^n)=n$, this proves the equivalence.
An automorphism preserves $\mathfrak m$ and all its powers. Applying
the equivalence to $a$ and its image shows that their valuations have
the same lower bounds in $\mathbb N$. Such lower bounds determine an
element of $\mathbb N\cup\{+\infty\}$, proving invariance.
\end{proof}

\begin{lemma}[The integral order function]
\label{fw:integral-order-function}
\dcref{Serre1979, Chapter IV, Section 1, Lemma 1 and Proposition 1}
\source{serre-1979-local-fields:page-0068,serre-1979-local-fields:page-0069}
\uses{fw:integer-lower-groups,fw:ideal-power-valuation}
\lean{FontaineWintenberger.ramificationOrder, FontaineWintenberger.ramificationOrder_le, FontaineWintenberger.mem_lowerRamificationGroup_iff_le_order, FontaineWintenberger.ramificationOrder_one, FontaineWintenberger.ramificationOrder_inv, FontaineWintenberger.ramificationOrder_conj, FontaineWintenberger.min_order_le_order_mul, FontaineWintenberger.ramificationOrder_mul_of_ne, FontaineWintenberger.ramificationOrder_eq_top_iff, FontaineWintenberger.lowerRamificationGroup_eventually_eq_bot, ValuationSubring.ramificationOrder, ValuationSubring.mem_lowerRamificationGroup_iff_le_order, ValuationSubring.ramificationOrder_eq_top_iff, ValuationSubring.lowerRamificationGroup_eventually_eq_bot}
\leanok
Let $D$ act by ring automorphisms on a discrete valuation ring $A$.
Define $j(s)=\inf_{a\in A}v_A(s(a)-a)$, with values in
$\mathbb N\cup\{+\infty\}$. Then
\[
 s\in D_n\ \Longleftrightarrow\ \max(n+1,0)\leq j(s),\qquad
 j(1)=+\infty,
\]
\[
 j(s^{-1})=j(s),\quad j(rsr^{-1})=j(s),\quad
 j(st)\geq\min\{j(s),j(t)\},
\]
and the last inequality is an equality when $j(s)\ne j(t)$.
If the action is faithful, $j(s)=+\infty$ exactly when $s=1$.
If in addition $D$ is finite, $D_n=\{1\}$ for all sufficiently large $n$.
In particular these assertions apply to the decomposition group of a
valuation subring that is a discrete valuation ring, and the last
conclusion holds if the ambient field has only finitely many base-field
automorphisms.
\end{lemma}
\begin{proof}
The membership equivalence follows by applying the ideal-power valuation
criterion to every difference and then taking its infimum.
The identity assertion follows from $v_A(0)=+\infty$.
Applying $s$ to $s^{-1}(a)-a$ gives $a-s(a)$; valuation invariance
and invariance under negation prove the inverse identity.
The difference $rsr^{-1}(a)-a$ is
$r(s(r^{-1}(a))-r^{-1}(a))$. Valuation invariance and the bijection
$a\mapsto r^{-1}(a)$ prove conjugation invariance.
The identity
\[
 st(a)-a=(s(t(a))-t(a))+(t(a)-a)
\]
and the valuation inequality for addition give the product inequality.
If $j(s)<j(t)$ and $j(st)>j(s)$, apply that inequality to
$s=(st)t^{-1}$ to contradict the strict inequalities.
If $j(t)<j(s)$, apply it to $t=s^{-1}(st)$ instead.
The infimum is infinite exactly when every difference vanishes;
faithfulness makes this equivalent to $s=1$.
For a finite faithful action, the finitely many values $j(s)$ for
$s\ne1$ are finite. An integer larger than all of them excludes every
nonidentity element from all subsequent groups.
\end{proof}

\begin{lemma}[Depth raising and the wild ramification group]
\label{fw:wild-ramification-p-group}
\source{serre-1979-local-fields:page-0073,serre-1979-local-fields:page-0074,serre-1979-local-fields:page-0075}
\dcref{Serre, Chapter IV, Section 2, Proposition 7 and Corollary 3}
\uses{fw:integer-lower-groups,fw:ideal-power-valuation,fw:integral-order-function}
\lean{FontaineWintenberger.smul_sub_mem_maximalIdeal_pow_add, FontaineWintenberger.pow_mem_lowerRamificationGroup_succ, FontaineWintenberger.pow_pow_mem_lowerRamificationGroup, FontaineWintenberger.lowerRamificationGroup_one_isPGroup}
\leanok
Let a finite group $D$ act faithfully by ring automorphisms on a
DVR $A$ with residue field of characteristic $p>0$. Write
$\mathfrak m$ for its maximal ideal and $D_n$ for the actual
integer congruence groups. If $s\in D_n$, $n\geq0$, and
$x\in\mathfrak m^k$, $k\geq1$, then
\[
 s(x)-x\in\mathfrak m^{n+k}.
\]
For $n\geq1$, $s\in D_n$ implies $s^p\in D_{n+1}$.
Consequently $D_1$ is a $p$-group. No finiteness or perfection
assumption on the residue field is required for this assertion.
\end{lemma}
\begin{proof}
Choose a uniformizer $\pi$ and write $x=\pi^k b$.
Both $\pi$ and $s(\pi)$ lie in $\mathfrak m$, and
$s(\pi)-\pi\in\mathfrak m^{n+1}$. The factorization
\[
 s(\pi)^k-\pi^k=(s(\pi)-\pi)
   \sum_{j=0}^{k-1}s(\pi)^{k-1-j}\pi^j
\]
puts this difference in $\mathfrak m^{n+k}$.
Now
\[
 s(\pi^k b)-\pi^k b=
   (s(\pi)^k-\pi^k)s(b)+\pi^k(s(b)-b).
\]
The first summand has the required order, and the second lies
in $\mathfrak m^{n+k+1}\subseteq\mathfrak m^{n+k}$.
This proves the first assertion, including $x=0$.

Let $n\geq1$ and put $d=s(a)-a\in\mathfrak m^{n+1}$.
Every power $s^j$ belongs to $D_n$. Applying the first assertion
to $s^j$ and $d$ gives
$s^j(d)-d\in\mathfrak m^{2n+1}\subseteq\mathfrak m^{n+2}$.
The telescoping identity
\[
 s^p(a)-a=\sum_{j=0}^{p-1}s^j(d)
\]
therefore makes this difference congruent to $pd$ modulo
$\mathfrak m^{n+2}$. Since the residue characteristic is $p$,
$p\in\mathfrak m$, and $pd\in\mathfrak m^{n+2}$.
Thus $s^p\in D_{n+1}$ by its actual congruence definition.

Induction now gives $s^{p^j}\in D_{j+1}$ for every $s\in D_1$.
Finite faithful actions have $D_N=1$ for some $N$; monotonicity
then gives $s^{p^N}=1$. Every element of the finite group $D_1$
has prime-power order, which is the assertion that $D_1$ is a
$p$-group. This proves the p-group conclusion of the cited
graded-quotient argument directly through depth raising, without
requiring the full residue-unit quotient construction first.
\end{proof}

\begin{lemma}[The uniformizer criterion over a perfect residue field]
\label{fw:perfect-residue-uniformizer-criterion}
\source{serre-1979-local-fields:page-0073}
\dcref{Serre, Chapter IV, Section 2, Proposition 5}
\uses{fw:integer-lower-groups,fw:ideal-power-valuation}
\lean{FontaineWintenberger.mem_lowerRamificationGroup_succ_of_uniformizer, FontaineWintenberger.mem_lowerRamificationGroup_iff_uniformizer}
\leanok
Let a group $D$ act by ring automorphisms on a DVR $A$ with
perfect residue field $k$ of characteristic $p>0$, and let $\pi$
be a uniformizer. For $s\in D_0$ and every integer $m\geq0$,
\[
 s\in D_m\quad\Longleftrightarrow\quad
 s(\pi)-\pi\in\mathfrak m^{m+1}.
\]
More precisely, if $s\in D_n$, $n\geq0$, and
$s(\pi)-\pi\in\mathfrak m^{n+2}$, then $s\in D_{n+1}$.
No completeness or finite-residue hypothesis is needed.
\end{lemma}
\begin{proof}
We first prove the more precise assertion. For each $x\in A$,
perfection of $k$ supplies $b\in A$ whose residue satisfies
$\overline b^p=\overline x$. Since $\mathfrak m=(\pi)$,
write $x=b^p+\pi c$ with $c\in A$.
Put $d=s(b)-b\in\mathfrak m^{n+1}$ and work modulo
$I=\mathfrak m^{n+2}$. Both $d^2$ and $pd$ vanish there:
$2(n+1)\geq n+2$ and $p\in\mathfrak m$.
The polynomial Taylor identity modulo a square-zero increment gives
\[
 (b+d)^p=b^p+pb^{p-1}d=b^p\pmod I.
\]
Thus $s(b^p)-b^p\in I$. Also
\[
 s(\pi c)-\pi c=(s(\pi)-\pi)s(c)+\pi(s(c)-c)\in I,
\]
using the extra uniformizer hypothesis and $s\in D_n$.
Adding proves $s(x)-x\in I$ for every $x$, as required.

The forward implication of the criterion follows directly from the
congruence definition. For the reverse implication use induction on
$m$. At $m=0$ the given hypothesis $s\in D_0$ already suffices.
If the uniformizer difference belongs to $\mathfrak m^{m+2}$,
it also belongs to $\mathfrak m^{m+1}$, so induction gives
$s\in D_m$, and the more precise assertion gives $s\in D_{m+1}$.
This proves the criterion. The argument adapts the cited criterion
by using perfect-residue Frobenius lifts in place of first constructing
the maximal unramified subfield.
\end{proof}

\begin{lemma}[The actual tame inertia quotient]
\label{fw:tame-inertia-quotient}
\source{serre-1979-local-fields:page-0073,serre-1979-local-fields:page-0074}
\dcref{Serre, Chapter IV, Section 2, Proposition 7 and Corollary 1}
\uses{fw:perfect-residue-uniformizer-criterion,fw:integer-lower-groups}
\lean{FontaineWintenberger.uniformizerRatio, FontaineWintenberger.mul_uniformizerRatio, FontaineWintenberger.uniformizerRatio_one, FontaineWintenberger.coe_uniformizerRatio_mul, FontaineWintenberger.residue_smul_eq_of_mem_lower_zero, FontaineWintenberger.tameInertiaHom, FontaineWintenberger.tameInertiaHom_eq, FontaineWintenberger.tameInertiaHom_coe, FontaineWintenberger.ker_tameInertiaHom, FontaineWintenberger.tameInertiaQuotientEquiv, FontaineWintenberger.isCyclic_tameInertiaQuotient, FontaineWintenberger.coprime_card_tameInertiaQuotient}
\leanok
Under the preceding hypotheses let $D$ be finite.
For $s\in D_0$ define
\[
 \theta(s)=\overline{s(\pi)/\pi}\in k^\times.
\]
The ratio is an actual unit of $A$. The map $\theta$ is a group
homomorphism, independent of $\pi$, with kernel $D_1$.
It identifies the actual quotient $D_0/D_1$ with its image in
$k^\times$. In particular $D_0/D_1$ is cyclic and its order is
prime to $p$.
\end{lemma}
\begin{proof}
Automorphisms preserve the maximal ideal and its generator order,
so $s(\pi)=\pi u_s$ for a unique unit $u_s$. Direct multiplication
gives $u_{st}=u_s s(u_t)$. Since $s\in D_0$ acts trivially on
residues, $\overline{s(u_t)}=\overline{u_t}$, proving
multiplicativity of $\theta$. If $\pi'=\pi v$ for a unit $v$,
then $s(\pi')/\pi'=u_s s(v)/v$; the second factor reduces to one,
proving independence of the uniformizer.
The condition $\theta(s)=1$ is equivalent to
$u_s-1\in\mathfrak m$, hence to
$s(\pi)-\pi\in\mathfrak m^2$. The uniformizer criterion
identifies this with $s\in D_1$. The first isomorphism theorem
therefore gives the stated quotient identification.

Every finite subgroup $H$ of the multiplicative group of a field
is cyclic. Explicitly, let $e$ be its exponent. An element of order
$e$ exists in the finite abelian group $H$: for each prime dividing
$e$, take an element with maximal prime-power order and multiply
these commuting elements of pairwise coprime orders. Every element
of $H$ is a root of $X^e-1$, so $|H|\leq e$ by the root bound;
the element of order $e$ gives the reverse inequality. Thus it
generates $H$. In characteristic $p$, the equation $z^p=1$
forces $(z-1)^p=0$ and hence $z=1$. A cyclic group of order
divisible by $p$ has an element of order $p$, so $p$ cannot divide
$|H|$. Apply this to the actual finite image of $\theta$.
\end{proof}

\begin{lemma}[Relative prime-step chains in a finite p-group]
\label{fw:relative-prime-subgroup-chain}
\source{serre-1979-local-fields:page-0075,wintenberger:page-0010}
\uses{fw:wild-ramification-p-group}
\lean{FontaineWintenberger.exists_normal_index_prime_le, FontaineWintenberger.subgroup_prime_chain_of_isPGroup, FontaineWintenberger.subgroup_prime_chain, FontaineWintenberger.wild_subgroup_prime_chain}
\leanok
Let $D$ be a finite group and $H\subseteq K\subseteq D$ actual
subgroups with $K$ a $p$-group. There is a finite descending chain
\[
 K=K_0\supseteq K_1\supseteq\cdots\supseteq K_d=H
\]
with $K_{j+1}$ normal in $K_j$ and $[K_j:K_{j+1}]=p$
for every $j<d$. The chain is allowed to be empty when $H=K$.
In particular this applies to every subgroup of the actual wild
ramification group of a finite faithful DVR action in residue
characteristic $p$.
\end{lemma}
\begin{proof}
We recall why maximal proper subgroups of a finite $p$-group $Q$
are normal of index $p$. A nontrivial finite $p$-group has
nontrivial center: in its conjugacy class equation every noncentral
class has cardinality divisible by $p$, so the cardinality of the
center is divisible by $p$; it is positive since it contains the
identity. Inducting on $|Q|$ after quotienting by the center shows
that $Q$ has a finite ascending central series. For any proper
subgroup $M$ choose the first term $Z_i$ of that series not contained
in $M$. Then $Z_{i-1}\subseteq M$, and an element of
$Z_i\setminus M$ normalizes $M$, since its commutators with elements
of $M$ lie in $Z_{i-1}$. Thus $M$ is properly contained in its
normalizer. If $M$ is maximal proper, its normalizer is all of $Q$,
so $M$ is normal. The quotient $Q/M$ is a nontrivial simple
$p$-group. Its nontrivial center is all of the quotient, so it is
abelian; any nonidentity element then generates it by simplicity,
and its cyclic prime-power order must be $p$ by simplicity.
Consequently $[Q:M]=p$.

Now induct on the finite cardinality of $K$. If $H=K$ take the
length-zero chain. Otherwise, among the proper subgroups of $K$
containing $H$ choose a maximal one $M$; it exists because $K$
is finite. The preceding paragraph gives $M\triangleleft K$ and
$[K:M]=p$. Since $|M|<|K|$, induction supplies a chain from
$M$ to $H$. Prepending $K$ gives the asserted chain. All subgroups
in this construction are subgroups of the original $D$, and each
normality and index assertion uses its actual inclusion into the
preceding term. The proved wild-ramification p-group theorem gives
the final assertion for the actual $D_1$. This is the group-theoretic
prime-layer refinement used in the cited norm reduction.
\end{proof}

\begin{lemma}[Actual relative cyclic prime-degree field towers]
\label{fw:relative-prime-field-tower}
\source{wintenberger:page-0010}
\uses{fw:relative-prime-subgroup-chain,fw:profinite-galois-foundations,fw:finite-extension-fixing-index,fw:wild-ramification-p-group}
\lean{FontaineWintenberger.isGalois_extendScalars_of_fixing_normal, FontaineWintenberger.cyclic_prime_extendScalars_of_fixing_step, FontaineWintenberger.relative_prime_field_chain, FontaineWintenberger.wild_relative_prime_field_chain}
\leanok
Let $L/K$ be finite Galois and let $E\subseteq M$ be actual
intermediate fields. If the actual fixing group $\Gal(L/E)$
is a $p$-group, there is a finite tower of actual intermediate fields
\[
 E=E_0\subseteq E_1\subseteq\cdots\subseteq E_d=M
\]
in which each $E_{j+1}/E_j$ is Galois, cyclic, and of degree $p$.
For a finite Galois extension of complete DVR fields in residue
characteristic $p$, this applies with $E$ the actual fixed field
of the first wild ramification group and every $M$ containing $E$.
\end{lemma}
\begin{proof}
Galois correspondence identifies the fixing groups of the actual
intermediate fields and reverses their inclusions. Apply the
relative prime-step subgroup chain to
$\Gal(L/M)\subseteq\Gal(L/E)$. Take the actual fixed fields of
its terms. The first and last are $E$ and $M$ by finite Galois
correspondence. Normality of each subgroup in the preceding one
implies normality, hence the Galois property, of the corresponding
field extension: restriction of the preceding fixing group onto
the smaller fixed field's automorphisms is surjective and has the
next subgroup as kernel. Its quotient is therefore that actual
relative Galois group. The subgroup index is the actual relative
field degree, so this quotient has order $p$. A group of prime
order is cyclic: any nonidentity element has order dividing $p$
and different from one, so generates the group. This proves every
assertion about each actual step.

In the valued-field case the actual first wild group is a finite
$p$-group by the proved wild ramification theorem. The fixing group
of its actual fixed field equals that group by Galois correspondence.
Thus the preceding argument applies with no assumed prime-step
tower or extra condition on $M/E$. The tower may have length zero
when $M=E$.
\end{proof}

\begin{lemma}[The normalized field valuation]
\label{fw:normalized-field-valuation}
\dcref{Serre1979, Chapter IV, Section 1, normalized-valuation convention}
\source{serre-1979-local-fields:page-0068}
\uses{fw:ideal-power-valuation,fw:integer-lower-groups}
\lean{FontaineWintenberger.mem_maximalIdeal_pow_iff_valuation, ValuationSubring.normalizedValuation_valuationSubring, ValuationSubring.mem_maximalIdeal_pow_iff_normalizedValuation, ValuationSubring.normalizedValuation_apply, ValuationSubring.mem_lowerRamificationGroup_iff_normalizedValuation}
\leanok
Let $A$ be a valuation subring of a field $E$, and assume $A$ is a
discrete valuation ring. The normalized multiplicative valuation
$w_A:E\to\{0\}\cup\{\exp(m):m\in\mathbb Z\}$ associated to its
maximal ideal has valuation ring exactly $A$. For $a\in A$ and
$n\geq0$,
\[
 a\in\mathfrak m_A^n\quad\Longleftrightarrow\quad
 w_A(a)\leq\exp(-n).
\]
Every $K$-automorphism of $E$ preserving $A$ preserves $w_A$.
Thus the integer congruence groups for the decomposition action are
exactly the groups defined by these normalized field-valuation bounds.
This supplies the standard fraction-field construction underlying the
normalized-valuation convention in the cited passage.
\end{lemma}
\begin{proof}
The field $E$ is the fraction field of $A$, since every field element or
its inverse belongs to $A$. Extend the valuation of $A$ to fractions by
$w_A(a/b)=w_A(a)/w_A(b)$. For nonzero $a,b$ this is
$\exp(v_A(b)-v_A(a))$. If $a/b=c/d$ with all four elements nonzero,
then $ad=bc$ and additivity gives
$v_A(a)+v_A(d)=v_A(b)+v_A(c)$, proving independence of the fraction.
The zero value is independent as well, since a fraction is zero exactly
when its numerator is zero. Multiplication follows from additivity,
and the inequality for addition follows by a common denominator and
the integral valuation inequality. Finally, $w_A(a/b)\leq1$ is
equivalent to $v_A(b)\leq v_A(a)$, hence to $b\mid a$, so the unit
ball is exactly $A$.
On a nonzero integral element its value is $\exp(-v_A(a))$;
at zero it is zero. The ideal-power criterion follows.
An automorphism preserving $A$ preserves its normalized integral
valuation. Therefore it preserves $w_A$ on $A$, and, by passing to
inverses, on every element of $E$. The final assertion is obtained by
applying the ideal-power criterion to every displacement.
\end{proof}

\begin{lemma}[A compatible norm for a discrete valuation]
\label{fw:discrete-compatible-norm}
\dcref{Serre1979, Chapter II, Section 2, Proposition 3 and its topology comparison}
\source{serre-1979-local-fields:page-0035,mathlib-normed-valued,mathlib-discrete-rank-one}
\lean{FontaineWintenberger.discreteValuedNormedField, FontaineWintenberger.discreteValued_norm_le_iff, FontaineWintenberger.discreteValued_valuationSubring, FontaineWintenberger.discreteValued_uniformSpace, FontaineWintenberger.discreteValued_completeSpace, FontaineWintenberger.discreteValued_unitBall_isDiscreteValuationRing}
\leanok
Let $K$ be a field with a nontrivial rank-one discrete valuation $v$,
endowed with its valuation uniformity. Its value group admits a strictly
increasing multiplicative embedding $\rho$ into $\mathbb R_{\geq0}$,
sending a positive generator to $2$. The norm $|x|=\rho(v(x))$
is nontrivial and nonarchimedean, defines the same uniformity, and satisfies
\[
 |x|\leq|y|\iff v(x)\leq v(y),\qquad
 |x|\leq1\iff x\in\mathcal O_v.
\]
Consequently its unit ball is exactly $\mathcal O_v$. If the valued field
is complete and $\mathcal O_v$ is a DVR, the resulting normed field is
complete with the same DVR as unit ball.
\end{lemma}
\begin{proof}
The nonzero value group is infinite cyclic with its ordered generator
chosen greater than one. Map its $n$th power to $2^n$, and send zero
to zero. This is multiplicative and strictly increasing. Composing
with $v$ gives multiplicativity, positive definiteness and the
ultrametric inequality, as well as the two displayed comparisons.
The positive value-group radii are cofinal among positive real radii:
powers of $2$ tend both to zero and to infinity. Thus the valuation
entourages $v(x-y)<\gamma$ and the norm entourages $|x-y|<\varepsilon$
define the same uniformity. Completeness is therefore preserved, and
the unit-ball equality transfers the DVR structure without changing
any element or ring operation.
\end{proof}

\begin{lemma}[Spectral generation over the source valuation ring]
\label{fw:source-spectral-generation}
\dcref{Serre1979, Chapter II, Sections 2 and 4; Chapter III, Section 6, Proposition 12}
\source{serre-1979-local-fields:page-0035,serre-1979-local-fields:page-0041,serre-1979-local-fields:page-0064,mathlib-normed-valued}
\uses{fw:discrete-compatible-norm,fw:adic-complete-unit-ball,fw:henselian-integral-generation,fw:source-chosen-extension-generation}
\lean{FontaineWintenberger.discreteValued_henselianLocalRing, FontaineWintenberger.discreteValued_spectral_hasExtension, FontaineWintenberger.discreteValued_spectral_exists_generator}
\leanok
Suppose the valued field in the compatible-norm lemma is complete and
its original valuation ring $R=\mathcal O_v$ is a DVR. Then $R$ is
Henselian. For an algebraic extension $E/K$, the spectral valuation
for the compatible norm extends $v$ in the sense of preserving the
ordering of values on $K$. If $E/K$ is finite separable and the residue
field of $R$ is perfect, its spectral integer ring is $R[x]$ for some
integral $x$. Here $R$ is the original source subring of $K$.
\end{lemma}
\begin{proof}
The compatible norm has unit ball exactly $R$ and preserves the source
uniformity and completeness. The adic completeness lemma therefore
makes $R$ Henselian. On $K$, the spectral norm equals the compatible
norm; its value comparisons are exactly those of $v$. This proves
extension compatibility. The spectral monogenicity theorem applies
to its norm unit ball. Identify that unit ball with $R$ by the identity
on elements of $K$. Polynomial evaluation commutes with this identity
on coefficients, so each polynomial representation of an integer
becomes a polynomial representation over the original $R$.
\end{proof}


\begin{theorem}[Extension of a nonarchimedean norm]
\label{fw:spectral-norm-extension}
\dcref{MathlibSpectralNorm, norm-extension theorem}
\source{mathlib-spectral-norm}
\lean{spectralMulAlgNorm, spectralNorm_extends, spectralNorm_unique_field_norm_ext, spectralNorm.completeSpace}
\mathlibok
Let $K$ be a complete nontrivially normed nonarchimedean field and
$E/K$ an algebraic extension. The spectral norm is a multiplicative
nonarchimedean norm extending the norm of $K$. Every absolute value on
$E$ extending that norm equals the spectral norm. If $E/K$ is finite,
$E$ is complete for the spectral norm.
\end{theorem}
\begin{proof}
These are the norm-extension and finite-dimensional completeness
theorems proved in \cite{MathlibSpectralNorm}. That source establishes
power-multiplicativity and the ultrametric inequality from the spectral
value, then uses completeness of the base to prove uniqueness and
multiplicativity. Its finite-dimensional completeness theorem applies
to the resulting normed vector space. This is the general algebraic
norm theorem; \cite[Chapter II, Section 2, Proposition 3]{Serre1979}
is its finite discretely valued specialization used here.
\end{proof}

\begin{lemma}[The spectral unit ball and integral closure]
\label{fw:spectral-integers}
\dcref{Serre1979, Chapter IV, notation and hypotheses}
\source{serre-1979-local-fields:page-0068}
\uses{fw:spectral-norm-extension}
\lean{FontaineWintenberger.isIntegral_iff_spectralNorm_le_one, FontaineWintenberger.spectralValuation, FontaineWintenberger.spectralValuation_algebraMap, FontaineWintenberger.spectralValuation_hasExtension, FontaineWintenberger.spectralIntegerRing, FontaineWintenberger.mem_spectralIntegerRing, FontaineWintenberger.spectralIntegerRing_eq_integralClosure, FontaineWintenberger.spectralIntegerRing_isIntegralClosure, FontaineWintenberger.smul_spectralIntegerRing, FontaineWintenberger.spectralIntegerRing_decompositionSubgroup}
\leanok
Let $K$ be a nonarchimedean normed field, let
$A_K=\{a\in K:|a|\leq1\}$, and let $E/K$ be algebraic.
Define the spectral norm of $x\in E$ to be the spectral value of its
monic minimal polynomial over $K$. Then
\[
 x\text{ is integral over }A_K\quad\Longleftrightarrow\quad
 |x|_{\mathrm{sp}}\leq1.
\]
If $K$ is complete and its norm is nontrivial, the spectral norm is a
multiplicative nonarchimedean norm extending the given norm, so its
unit ball is a valuation subring of $E$. This valuation subring is the
integral closure of $A_K$ and is preserved by every $K$-automorphism of $E$.
This identifies the integral-closure description in the cited passage
with the spectral-norm construction; the proof uses minimal-polynomial
coefficients and does not require a finite residue field.
\end{lemma}
\begin{proof}
The spectral value of a monic polynomial $P$ is at most one if and only
if every coefficient of $P$ has norm at most one. Indeed its defining
terms below the leading degree are
$|P_i|^{1/(\deg P-i)}$, and the exponents are positive; the leading
coefficient is one and coefficients above the degree vanish.
The valuation ring $A_K$ is integrally closed: if an element of $K$ of
norm greater than one satisfied a monic equation with coefficients of
norm at most one, its leading term would have strictly larger norm than
all other terms, contrary to that equation.
Consequently, for an integral $x$, its minimal polynomial over $K$ is
the image of its monic minimal polynomial over $A_K$. Its coefficients
belong to $A_K$, giving the forward implication. Conversely, if
$|x|_{\mathrm{sp}}\leq1$, all coefficients of its monic minimal
polynomial belong to $A_K$; viewed as a polynomial over $A_K$, it gives
the required monic equation for $x$.
For a complete nontrivially normed base, Theorem~\ref{fw:spectral-norm-extension}
makes the spectral norm multiplicative and nonarchimedean. Its unit
ball is therefore a valuation subring. The equivalence just proved
identifies this subring with the integral closure. Finally,
$K$-automorphisms preserve minimal polynomials, hence spectral norms
and their unit ball.
\end{proof}

\begin{lemma}[The full automorphism action on spectral integers]
\label{fw:spectral-full-galois-action}
\dcref{Serre1979, Chapter II, Section 2, Proposition 3; Chapter IV, Section 1}
\source{serre-1979-local-fields:page-0035,serre-1979-local-fields:page-0068}
\uses{fw:spectral-integers,fw:integer-lower-groups,fw:integral-order-function,fw:normalized-field-valuation}
\lean{FontaineWintenberger.spectralDecompositionEquiv, FontaineWintenberger.spectralIntegerRing_galoisAction, FontaineWintenberger.coe_spectralIntegerRing_galois_smul, FontaineWintenberger.spectralIntegerRing_galois_faithfulSMul, FontaineWintenberger.spectralIntegerRing_galois_smulCommClass, FontaineWintenberger.spectralIntegerRing_galois_mem_lower_iff, FontaineWintenberger.spectralIntegerRing_galois_order, FontaineWintenberger.spectralIntegerRing_galois_normalizedValuation}
\leanok
Let $E/K$ be an algebraic extension of a complete nontrivially normed
nonarchimedean field, and let $A_E$ be its spectral integer ring.
Its decomposition group equals $\operatorname{Aut}_K(E)$, with the
identification sending an automorphism to itself. Transport of the
intrinsic action gives a faithful action of this full group on $A_E$;
after inclusion into $E$, it is the original field action.
When $A_E$ is a discrete valuation ring, this identification preserves
every integer lower group, its intrinsic order function, and its
normalized valuation. The action fixes base integers.
\end{lemma}
\begin{proof}
Spectral norms are invariant under $K$-automorphisms, so every such
automorphism stabilizes $A_E$. Hence each belongs to the decomposition
group, whose inclusion is therefore a group isomorphism. Transport the
intrinsic action along its inverse. Inclusion into $E$ commutes with
this action by construction. Faithfulness follows from the faithful
intrinsic action. The base is fixed because the field automorphisms
are $K$-linear. The definitions of lower groups and order use exactly
the same integral displacements under this identification; normalized
valuation invariance follows from the intrinsic invariance lemma.
\end{proof}

\begin{lemma}[Finite extension valuation rings]
\label{fw:finite-valuation-ring}
\dcref{Serre1979, Chapter II, Section 2, Proposition 3, separable case}
\source{serre-1979-local-fields:page-0035,serre-1979-local-fields:page-0036}
\uses{fw:spectral-integers}
\lean{FontaineWintenberger.spectralIntegerRing_moduleFinite, FontaineWintenberger.spectralIntegerRing_moduleFree, FontaineWintenberger.spectralIntegerRing_finrank, FontaineWintenberger.spectralIntegerRing_isNoetherianRing, FontaineWintenberger.spectralIntegerRing_not_isField, FontaineWintenberger.spectralIntegerRing_isDiscreteValuationRing}
\leanok
Let $K$ be a complete nontrivially normed nonarchimedean field whose
unit ball $A_K$ is a discrete valuation ring. If $E/K$ is finite
separable, its spectral valuation ring $A_E$ is a discrete valuation
ring and a finite free $A_K$-module of rank $[E:K]$.
This proof uses the spectral unit ball to obtain the unique local ring,
in place of the comparison of prime-valuation topologies in Serre's proof.
\end{lemma}
\begin{proof}
By the spectral integrality criterion, $A_E$ is the integral closure
of $A_K$ in $E$. The trace-dual-basis bound for integral closure in a
finite separable extension makes $A_E$ a finite $A_K$-module, since
$A_K$ is Noetherian and integrally closed. It is torsion-free, being
a subring of $E$, and hence free over the principal ideal ring $A_K$.
Localizing at the nonzero elements of $A_K$ recovers $E$: for any
algebraic $x\in E$, multiplying by a sufficiently large power of a
base uniformizer makes its spectral norm at most one. Thus the rank
is $[E:K]$.
Finiteness also makes $A_E$ Noetherian. A Noetherian valuation ring is
a principal ideal ring. To see it is not a field, take a uniformizer
$\pi$ of $A_K$. Its image is nonzero and has spectral norm
$|\pi|<1$, so it is a nonunit of $A_E$. Therefore the local principal
ideal domain $A_E$ is a discrete valuation ring.
\end{proof}

\begin{lemma}[The spectral integers after changing a finite base]
\label{fw:spectral-base-integers}
\uses{fw:spectral-norm-extension,fw:spectral-integers,fw:normed-intermediate-tower}
\lean{FontaineWintenberger.spectralIntegerRing_eq_normedUnitBall, FontaineWintenberger.finiteBase_spectralIntegerRing_eq, FontaineWintenberger.finiteBase_spectralIntegerRingEquiv, FontaineWintenberger.coe_finiteBase_spectralIntegerRingEquiv, FontaineWintenberger.finiteBase_spectralIntegerRingEquiv_smul}
\leanok
Let $K$ be complete, nonarchimedean and nontrivially normed, let
$L/K$ be algebraic with a compatible nonarchimedean field norm,
and let $M$ be a finite intermediate field with its restricted norm.
The spectral valuation rings of $L$ constructed over $K$ and over
$M$ are the same actual subring of $L$, namely its norm unit ball.
Their identity isomorphism intertwines the action of the actual
fixing subgroup $\Gal(L/M)\subseteq\operatorname{Aut}_K(L)$
with its ordinary $M$-automorphism action. No normality or finite
degree of $L/K$ is required.
\end{lemma}
\begin{proof}
The finite intermediate field $M$ is complete in its restricted
norm, and $L$ has a compatible norm as an extension of each of
$K$ and $M$. For any complete nonarchimedean normed base, the
spectral norm on an algebraic normed extension equals its given
compatible norm. Thus both spectral unit balls are exactly
$\{x\in L:|x|\leq1\}$. Their identity as subrings gives the
claimed ring isomorphism. The fixing-subgroup equivalence sends
each automorphism to the same function on $L$. Both induced
actions on the equal integer subrings are its original action
on field elements, proving equivariance.
\end{proof}

\begin{lemma}[Spectral integers and residues along field embeddings]
\label{fw:spectral-integer-residue-maps}
\uses{fw:spectral-norm-extension,fw:spectral-integers}
\lean{FontaineWintenberger.spectralValuation_hasExtension_tower, FontaineWintenberger.spectralIntegerInclusion, FontaineWintenberger.coe_spectralIntegerInclusion, FontaineWintenberger.spectralIntegerInclusion_isLocalHom, FontaineWintenberger.spectralResidueInclusion, FontaineWintenberger.spectralResidueInclusion_residue, FontaineWintenberger.spectralResidueInclusion_injective, FontaineWintenberger.spectralIntegerInclusion_comp, FontaineWintenberger.spectralResidueInclusion_comp}
\leanok
Let $K$ be complete nonarchimedean and nontrivially normed, and let
$E,L,M$ be algebraic extensions of $K$. A $K$-embedding
$i:E\to L$ restricts to a local homomorphism of the actual
spectral valuation rings and induces an injective homomorphism
$k_E\to k_L$ of algebras over the original residue field $k_K$.
These integer and residue maps compose along the original field
embeddings. No finite-degree or discrete-valuation hypothesis on
the extensions is needed.
\end{lemma}
\begin{proof}
Use $i$ to regard $L$ as an extension of $E$. Exact restriction
of the spectral norm gives $|i(x)|_L=|x|_E$ for all $x\in E$.
It therefore sends the actual unit ball into the actual unit ball
and an element of norm less than one precisely to such an element.
The induced integer map is local, since units of a valuation ring
are exactly the elements whose element and inverse both lie in
that ring. It fixes the original base valuation ring. Passing to
the quotients by maximal ideals therefore gives a unital map of
fields, hence an injection, which fixes $k_K$. Its formula on
a residue class is the residue class of the image under $i$.
For composable embeddings this formula proves both composition
identities: every residue class has an integral representative,
and the underlying integer inclusions are the original embeddings.
\end{proof}

\begin{lemma}[Uniqueness of the extension valuation ring]
\label{fw:spectral-valuation-ring-unique}
\dcref{Serre1979, Chapter II, Section 2, Proposition 3}
\source{serre-1979-local-fields:page-0035,serre-1979-local-fields:page-0036}
\uses{fw:spectral-integers,fw:finite-valuation-ring}
\lean{FontaineWintenberger.spectralIntegerRing_eq_extension_valuationSubring}
\leanok
Let $K$ be complete, nontrivially normed and nonarchimedean, with DVR
unit ball $R$, and let $E/K$ be finite separable. Every valuation on
$E$ extending the valuation of $K$ up to equivalence has valuation ring
equal to the actual spectral integer ring. Thus this conclusion is
independent of the chosen target value group and normalization.
\end{lemma}
\begin{proof}
Write $B$ for the spectral integer ring and $W$ for the ring of another
compatible extension. Every element of $B$ is integral over $R$.
Since $R$ maps into $W$, it is also integral over $W$, and integral
closedness of the valuation ring $W$ puts it in $W$. Hence $B\subseteq W$.
The DVR $B$ has Krull dimension one, so its only valuation overrings
in its fraction field are $B$ and the field $E$: equivalently, overrings
correspond to the two prime ideals of $B$. The second case is impossible.
Indeed a base uniformizer has nonzero image which is a nonunit of $W$,
because the map $R\to W$ is local by valuation compatibility. A field
has no such element. Thus $W=B$.
\end{proof}

\begin{corollary}[Integers for a chosen compatible extension valuation]
\label{fw:chosen-extension-integers}
\dcref{Serre1979, Chapter II, Sections 2 and 4; Chapter III, Section 6, Proposition 12}
\source{serre-1979-local-fields:page-0035,serre-1979-local-fields:page-0036,serre-1979-local-fields:page-0041,serre-1979-local-fields:page-0064}
\uses{fw:spectral-valuation-ring-unique,fw:finite-valuation-ring,fw:finite-extension-henselian,fw:henselian-integral-generation}
\lean{Valuation.HasExtension.valuationSubring_faithfulSMul, FontaineWintenberger.extensionValuationSubring_isIntegralClosure, FontaineWintenberger.extensionValuationSubring_moduleFinite, FontaineWintenberger.extensionValuationSubring_isDiscreteValuationRing, FontaineWintenberger.extensionValuationSubring_henselianLocalRing, FontaineWintenberger.extensionValuationSubring_exists_generator}
\leanok
Under the preceding uniqueness hypotheses, the chosen valuation ring
$W$ is the actual integral closure of $R$ in $E$, is finite over $R$,
is a DVR and is Henselian. If the residue field of $R$ is perfect,
then $W=R[x]$ for some $x\in W$.
\end{corollary}
\begin{proof}
The ring equality identifies $W$ with the spectral integer ring by the
identity on elements of $E$, and all inclusion maps into $E$ agree.
It is therefore the integral closure, and integral-closure finiteness
applies to its actual $R$-algebra structure. Its DVR and Henselian
properties transfer along the same equality. The base-to-$W$ map is
injective, because its composite with $W\hookrightarrow E$ is the
inclusion from $R$ through $K$. Its residue extension is finite and
separable when the base residue field is perfect. Apply Henselian
integral generation to this actual finite faithful local algebra.
\end{proof}


\begin{lemma}[Adic completeness of the unit ball]
\label{fw:adic-complete-unit-ball}
\dcref{Serre1979, Chapter II, Section 4, Proposition 7}
\source{serre-1979-local-fields:page-0041}
\lean{FontaineWintenberger.mem_maximalIdeal_pow_iff_norm, FontaineWintenberger.unitBall_isAdic, FontaineWintenberger.unitBall_completeSpace, FontaineWintenberger.unitBall_isAdicComplete, FontaineWintenberger.unitBall_henselianLocalRing}
\leanok
Let $K$ be a nonarchimedean normed field whose unit ball $A$ is a discrete
valuation ring, and choose a uniformizer $\pi\in A$. Then
\[
 a\in\mathfrak m_A^n\quad\Longleftrightarrow\quad
 |a|\leq|\pi|^n.
\]
The norm topology of $A$ is its maximal-ideal-adic topology. If $K$ is
complete, then $A$ is complete for this adic topology and is Henselian.
\end{lemma}
\begin{proof}
The ideal power is $(\pi^n)$, and divisibility in a valuation ring is
equivalent to the reverse inequality of values, proving the displayed
criterion. Since $0<|\pi|<1$, these balls have positive radii tending
to zero. A positive-radius closed ball in an ultrametric space is open:
the open ball of the same radius about each of its points lies inside it.
Thus the powers are open and form a basis of neighborhoods of zero,
which identifies the two ring topologies. The set $A$ is closed in $K$,
being the closed unit ball. Completeness of $K$ therefore makes $A$
complete and Hausdorff. Completeness for the identified adic topology
gives Hensel lifting as follows. Start at a lift $a_0$ of a simple
residue root of $f$, and put
$a_{r+1}=a_r-f(a_r)/f'(a_r)$. The derivative is a unit throughout,
because all $a_r$ have the same residue. Taylor's formula shows
$f(a_r)\in\mathfrak m_A^{2^r}$ and
$a_{r+1}-a_r\in\mathfrak m_A^{2^r}$ by induction. Hence $(a_r)$ is
Cauchy. Its limit has the prescribed residue and is a root, by
continuity of polynomial evaluation. If $a,b$ are roots with this
residue, divide $f(T)$ by $T-a$ to write $f(T)=(T-a)g(T)$.
The reduction of $g(b)$ is the nonzero derivative at the residue root,
so $g(b)$ is a unit. Thus $(b-a)g(b)=0$ forces $b=a$.
\end{proof}

\begin{corollary}[Hensel lifting in a finite extension]
\label{fw:finite-extension-henselian}
\dcref{Serre1979, Chapter II, Sections 2 and 4, Propositions 3 and 7}
\source{serre-1979-local-fields:page-0035,serre-1979-local-fields:page-0041}
\uses{fw:spectral-norm-extension,fw:finite-valuation-ring,fw:adic-complete-unit-ball}
\lean{FontaineWintenberger.spectralIntegerRing_henselianLocalRing}
\leanok
Under the hypotheses of the finite extension valuation-ring lemma,
$A_E$ is Henselian.
\end{corollary}
\begin{proof}
The spectral norm makes $E$ complete. Its unit ball is precisely $A_E$,
which is a discrete valuation ring by the finite extension lemma.
The preceding result applies to this norm and unit ball.
\end{proof}

\begin{lemma}[The canonical ramification index and normalized valuation]
\label{fw:ramification-index-scaling}
\dcref{Serre1979, Chapter II, Section 2; Chapter III, Section 6, Lemma 3}
\source{serre-1979-local-fields:page-0035,serre-1979-local-fields:page-0064}
\uses{fw:ideal-power-valuation}
\lean{FontaineWintenberger.exists_uniformizer_factor_ramificationIdx, FontaineWintenberger.map_maximalIdeal_eq_pow_ramificationIdx, FontaineWintenberger.local_ramificationIdx_pos, FontaineWintenberger.ramificationIdx_eq_addVal_uniformizer, FontaineWintenberger.addVal_algebraMap}
\leanok
This is an algebraic DVR generalization of the cited finite-extension
facts; neither finiteness nor completeness is needed for this statement.
Let $R\hookrightarrow S$ be an injective local map of DVRs. Put
$e=\operatorname{ramificationIdx}_R(\mathfrak m_S)$, using the canonical
local-length ramification index. Then $e>0$,
\[
 \mathfrak m_RS=\mathfrak m_S^e,\qquad
 v_S(a)=e\,v_R(a)\quad(a\in R).
\]
Here the integral valuations take values in $\mathbb N\cup\{\infty\}$;
since $e>0$, the formula includes $a=0$.
\end{lemma}
\begin{proof}
Choose uniformizers $\varpi$ of $R$ and $\pi$ of $S$. Write the nonzero
image of $\varpi$ as $u\pi^d$, with $u\in S^\times$. Then
$\mathfrak m_RS=(\varpi)S=(\pi^d)=\mathfrak m_S^d$. This ideal is
not contained in $\mathfrak m_S^{d+1}$, by the valuation of $\pi^d$.
Thus its largest power exponent is $d$. The ideal-power definition
agrees with the canonical local-length index in the Dedekind domain
$S$, proving $e=d$. Localness makes the mapped maximal ideal proper,
so $d\ne0$. For $a\ne0$, write $a=v\varpi^n$ with $v\in R^\times$.
Its image has valuation $dn=en$, since units map to units. At zero
both valuations are infinite and $e$ is positive.
\end{proof}

\begin{theorem}[Adic valuation scaling over a prime]
\label{fw:adic-valuation-extension}
\source{mathlib-ramification-valuation}
\cite{MathlibRamificationValuation}
\lean{IsDedekindDomain.HeightOneSpectrum.intValuation_liesOver, IsDedekindDomain.HeightOneSpectrum.valuation_liesOver}
\mathlibok
Let $A\to B$ be an algebra of Dedekind domains with $B$ torsion-free
as an $A$-module. Let $\mathfrak q$ be a height-one prime of $B$ lying
over a height-one prime $\mathfrak p$ of $A$. The normalized adic values
satisfy
\[
 w_{\mathfrak q}(a)=w_{\mathfrak p}(a)^{e(\mathfrak q/\mathfrak p)}
 \qquad(a\in A).
\]
For compatible fraction fields the same formula holds on the lower
fraction field. This is the pinned imported adic-valuation theorem.
\end{theorem}
\begin{proof}
For nonzero $a$, the exponent of $\mathfrak q$ in $(a)B$ is the exponent
of $\mathfrak p$ in $(a)$ multiplied by the ramification index, by
unique ideal factorization and lying-over. The normalized values are
the corresponding negative powers of their fixed value-group generator,
so this exponent identity gives the displayed formula. At zero use
positivity of the ramification index. Passing to a quotient of two
integral elements proves the fraction-field assertion.
\end{proof}

\begin{lemma}[Normalized valuations and ramification in towers]
\label{fw:normalized-valuation-scaling}
\dcref{Serre1979, Chapter II, Section 2, Proposition 3}
\source{serre-1979-local-fields:page-0035,serre-1979-local-fields:page-0064}
\uses{fw:ramification-index-scaling,fw:normalized-field-valuation,fw:adic-valuation-extension,fw:chosen-extension-integers}
\lean{FontaineWintenberger.normalizedValuation_algebraMap, FontaineWintenberger.normalizedValuation_hasExtension, FontaineWintenberger.local_ramificationIdx_tower, FontaineWintenberger.extensionValuationSubring_map_maximalIdeal, FontaineWintenberger.extensionValuationSubring_normalizedValuation_algebraMap}
\leanok
Let $R\hookrightarrow S$ be an injective local map of DVRs, with
compatible fraction fields $K\hookrightarrow E$. For their canonical
normalized multiplicative valuations,
\[
 w_S(a)=w_R(a)^{e(S/R)}\qquad(a\in K).
\]
In particular the normalized valuation of $E$ extends the valuation of
$K$ up to equivalence. For a tower of injective local DVR maps
$R\hookrightarrow S\hookrightarrow T$,
\[
 e(T/R)=e(S/R)e(T/S).
\]
These formulas apply, in particular, to the actual base valuation ring
and every chosen compatible valuation ring in the finite-extension
setting above; the latter DVR structure is supplied by the chosen-ring
theorem.
\end{lemma}
\begin{proof}
On nonzero integral elements, factor into a unit times a power of a
uniformizer. Units retain value one, while the mapped base uniformizer
has integral order $e(S/R)$ by the preceding lemma. This gives the
power formula on $R$, also at zero since the index is positive.
Represent any element of $K$ as a quotient of elements of $R$ and use
multiplicativity to obtain the formula on $K$. Taking a positive
power preserves and reflects comparisons of nonnegative values, giving
extension equivalence. In a tower, map the equality
$\mathfrak m_RS=\mathfrak m_S^{e(S/R)}$ to $T$, and then use the analogous
equality for $S\to T$. Powers of a DVR maximal ideal have distinct
exponents, so their exponents multiply as claimed.
The chosen-extension assertions are these statements applied to the
actual fraction fields and the DVR structure already proved for the
chosen ring. Their canonical inclusion maps satisfy the required towers.
\end{proof}

\begin{lemma}[Generation from the residue field and a uniformizer]
\label{fw:uniformizer-residue-generation}
\dcref{Serre1979, Chapter III, Section 6, Proposition 12}
\source{serre-1979-local-fields:page-0064}
\lean{FontaineWintenberger.subalgebra_eq_top_of_residue_surjective}
\leanok
Let $R\hookrightarrow S$ be an injective map of discrete valuation rings,
with $S$ finite over $R$. If an $R$-subalgebra $A\subseteq S$ maps onto
$S/\mathfrak m_S$ and contains a uniformizer $z$ of $S$, then $A=S$.
\end{lemma}
\begin{proof}
Successive residue lifts show, for every $b\in S$ and $n\geq0$, that
$b=a+z^n c$ with $a\in A$ and $c\in S$. For $n=0$ take $a=0$.
Given a decomposition at $n$, choose $d\in A$ with the residue of $c$.
Then $c-d=zc'$ since $\mathfrak m_S=(z)$, and replace $a$ by
$a+z^nd\in A$ to obtain the decomposition at $n+1$.
Take a uniformizer $\varpi$ of $R$. Its nonzero image has the form
$uz^e$ for a unit $u\in S^\times$. The decomposition at $e$ therefore
gives $S=A+\mathfrak m_RS$. Apply Nakayama's lemma to the finite
$R$-module $S/A$: since $\mathfrak m_R$ is in the Jacobson radical,
this quotient vanishes.
\end{proof}

\begin{lemma}[Henselian integral generation]
\label{fw:henselian-integral-generation}
\dcref{Serre1979, Chapter III, Section 6, Proposition 12}
\source{serre-1979-local-fields:page-0064}
\uses{fw:uniformizer-residue-generation,fw:finite-valuation-ring,fw:finite-extension-henselian}
\lean{FontaineWintenberger.exists_monic_root_of_residue, FontaineWintenberger.irreducible_aeval_add_of_isRoot, FontaineWintenberger.exists_integral_generator, FontaineWintenberger.adjoin_uniformizer_eq_top, FontaineWintenberger.spectralIntegerRing_exists_generator}
\leanok
Let $R\hookrightarrow S$ be a finite injective local map of discrete
valuation rings. Suppose $S$ is Henselian and the residue extension is
separable. Then $S=R[x]$ for some $x\in S$.
If the map on residue fields is surjective, every uniformizer of $S$
generates $S$ over $R$, without a Henselian hypothesis.
In particular the spectral integer ring of a finite separable extension
of a complete nonarchimedean normed field with discrete valuation ring
as unit ball and perfect residue field is monogenic over that unit ball.
\end{lemma}
\begin{proof}
The residue extension is finite because $S$ is finite over $R$.
Choose a primitive residue element $\alpha$ and lift its monic minimal
polynomial to $f\in R[T]$. Separability says the derivative at $\alpha$
is nonzero. Hensel lifting in $S$ supplies a root $a$ of $f$ with residue
$\alpha$ and $f'(a)$ a unit. Choose a uniformizer $\pi$ of $S$ and put
$x=a+\pi$. Taylor expansion gives
$f(x)=\pi(f'(a)+c\pi)$ for some $c\in S$. The second factor is a unit
because its residue equals that of $f'(a)$. Thus $f(x)$ is a uniformizer
in $R[x]$, and the residue of $x$ generates the residue field. The
preceding lemma gives $R[x]=S$.
If the residue map is surjective, $R[\pi]$ already maps onto the residue
field and contains a uniformizer, so the same lemma applies directly.
For the stated spectral specialization, finite freeness, the DVR
structure and Henselianity are supplied by the preceding lemmas;
perfection of the base residue field makes its finite extension separable.
This proves the monogenic conclusion needed here, without asserting
Serre's additional power-basis conclusion as a separate new result.
\end{proof}

\begin{lemma}[Integral generation for source-valued finite extensions]
\label{fw:source-chosen-extension-generation}
\dcref{Serre1979, Chapter II, Section 2, Proposition 3; Chapter III, Section 6, Proposition 12}
\source{serre-1979-local-fields:page-0035,serre-1979-local-fields:page-0064,mathlib-normed-valued}
\uses{fw:discrete-compatible-norm,fw:chosen-extension-integers}
\lean{FontaineWintenberger.discreteValued_norm_hasExtension, FontaineWintenberger.discreteValued_extension_isDiscreteValuationRing, FontaineWintenberger.discreteValued_extension_henselianLocalRing, FontaineWintenberger.discreteValued_extension_exists_generator}
\leanok
Let $K$ be complete for a rank-one discrete valuation $v$, whose
valuation ring $R$ is a DVR with perfect residue field. For a finite
separable extension $E/K$ and any chosen valuation $w$ extending $v$
up to equivalence, its actual valuation ring $W$ is a DVR, is Henselian,
and satisfies $W=R[x]$ for some $x\in W$.
\end{lemma}
\begin{proof}
Choose the compatible norm of the discrete-compatible-norm lemma.
The norm-induced valuation has exactly the comparisons of $v$;
transitivity of equivalence therefore makes $w$ a compatible extension
of this norm-induced valuation. The chosen-extension theorem now
makes the actual $W$ a DVR and Henselian, and produces a monogenic
generator over the norm unit ball. That unit ball is exactly the
original source ring $R$. The identity on its elements commutes with
the inclusions into $E$ and $W$. Mapping polynomial coefficients by
this identity transports every polynomial representation, giving
$W=R[x]$ with the original source algebra structure.
\end{proof}

\begin{theorem}[The finite-flat degree identity]
\label{fw:finite-flat-degree-identity}
\dcref{MathlibRamificationInertia, finite-flat degree identity and maximal-ideal inertia degree}
\source{mathlib-ramification-inertia-basic,mathlib-inertia-degree}
\lean{Ideal.sum_ramification_inertia_eq_finrank,Ideal.inertiaDeg_eq_of_isMaximal}
\mathlibok
Let $R$ be a domain and $S$ a finite flat commutative $R$-algebra.
For a prime ideal $\mathfrak p$ of $R$ with finitely many primes above it,
\[
 \sum_{\mathfrak q\mid\mathfrak p} e(\mathfrak q/\mathfrak p)
 f(\mathfrak q/\mathfrak p)=\operatorname{rank}_R S.
\]
For maximal $\mathfrak p,\mathfrak q$ with $\mathfrak q$ over $\mathfrak p$,
the canonical inertia degree equals
$[S/\mathfrak q:R/\mathfrak p]$.
\end{theorem}
\begin{proof}
We use the local-length argument of \cite{MathlibRamificationInertia}.
Put $A=S\otimes_R\kappa(\mathfrak p)$. The finite flat module
$S_{\mathfrak p}$ is free over the local ring $R_{\mathfrak p}$.
Its rank equals the dimension after tensoring further with
$\operatorname{Frac}(R)$, hence equals $\operatorname{rank}_R S$.
Reducing this free module modulo $\mathfrak p$ gives
$\dim_{\kappa(\mathfrak p)} A=\operatorname{rank}_R S$.
The finite-dimensional algebra $A$ is Artinian and is the product of
its localizations at its finitely many maximal ideals. These ideals
correspond to the primes $\mathfrak q$ above $\mathfrak p$.
The corresponding local factor is
$S_{\mathfrak q}/\mathfrak pS_{\mathfrak q}$. A composition series of
this local Artinian ring as a module over itself has
$e(\mathfrak q/\mathfrak p)$ factors, by the definition of the
ramification index, and each simple factor is
$\kappa(\mathfrak q)$. Each such factor has
$\kappa(\mathfrak p)$-dimension $f(\mathfrak q/\mathfrak p)$.
Additivity of dimension along this series and then along the finite
product gives the asserted sum. If both ideals are maximal,
$\kappa(\mathfrak p)$ and $\kappa(\mathfrak q)$ are respectively the
quotient fields $R/\mathfrak p$ and $S/\mathfrak q$, and their canonical
structure map is the quotient of $R\to S$, proving the last assertion.
\end{proof}

\begin{lemma}[The degree identity for local DVR extensions]
\label{fw:local-dvr-degree-identity}
\dcref{Serre1979, Chapter II, Section 2, Corollary 1 to Proposition 3}
\source{serre-1979-local-fields:page-0036}
\uses{fw:finite-flat-degree-identity,fw:chosen-extension-integers}
\lean{FontaineWintenberger.local_inertiaDeg_eq_finrank_residueField, FontaineWintenberger.local_ramificationIdx_mul_inertiaDeg, FontaineWintenberger.local_ramificationIdx_mul_finrank_residueField, FontaineWintenberger.local_ramificationIdx_mul_residueDegree_eq_finrank, FontaineWintenberger.extensionValuationSubring_ramificationIdx_mul_residueDegree, FontaineWintenberger.discreteValued_extension_moduleFinite, FontaineWintenberger.discreteValued_extension_ramificationIdx_mul_residueDegree}
\leanok
Let $R\hookrightarrow S$ be a finite injective local map of DVRs.
Write $e=e(S/R)$ for the canonical ramification index and
$f=[k_S:k_R]$ for the residue-field degree. Then
\[
 ef=\operatorname{rank}_R S.
\]
If $K$ and $E$ are compatible fraction fields of $R$ and $S$, this
rank is $[E:K]$. In particular these assertions apply to an actual
chosen extension valuation ring of a finite separable extension of a
complete nonarchimedean field with DVR unit ball.
\end{lemma}
\begin{proof}
The injective map makes $S$ torsion-free over the PID $R$, hence free
and flat. Since $S$ is finite over $R$, it is integral over $R$.
Every prime over the maximal ideal of $R$ is therefore maximal, and
the local ring $S$ has exactly one such prime, its maximal ideal.
The finite-flat identity thus has one summand, giving $ef$ equal
to the rank. Localization of a finite free module preserves its
rank; localizing $R$ to $K$ and $S$ to its fraction field identifies
this rank with $[E:K]$. For the last assertion use the already proved
finiteness and DVR structure of the actual chosen valuation ring.
\end{proof}

\begin{lemma}[Total ramification and every uniformizer]
\label{fw:totally-ramified-uniformizer-generation}
\dcref{Serre1979, Chapter II, Section 2, Corollary 1; Chapter III, Section 6, Proposition 12}
\source{serre-1979-local-fields:page-0036,serre-1979-local-fields:page-0064}
\uses{fw:local-dvr-degree-identity,fw:ramification-index-scaling,fw:henselian-integral-generation,fw:source-chosen-extension-generation}
\lean{FontaineWintenberger.ramificationIdx_eq_finrank_iff_residue_surjective, FontaineWintenberger.adjoin_uniformizer_eq_top_of_ramificationIdx_eq_finrank, FontaineWintenberger.adjoin_uniformizer_eq_top_of_ramificationIdx_eq_field_finrank, FontaineWintenberger.extensionValuationSubring_adjoin_uniformizer_eq_top, FontaineWintenberger.discreteValued_extension_adjoin_uniformizer_eq_top}
\leanok
For a finite injective local map of DVRs $R\hookrightarrow S$,
\[
 e(S/R)=\operatorname{rank}_R S
 \quad\Longleftrightarrow\quad k_R\longrightarrow k_S\text{ is surjective}.
\]
Under these equivalent conditions, every uniformizer $\pi$ of $S$
satisfies $R[\pi]=S$. Consequently, for a finite separable extension
$E/K$ of a complete rank-one discretely valued field with DVR
valuation ring, if $e(E/K)=[E:K]$, every uniformizer in its actual
chosen compatible extension valuation ring generates that ring over
the original source valuation ring. No perfection hypothesis is
needed for this totally ramified assertion.
\end{lemma}
\begin{proof}
The degree identity gives $ef=\operatorname{rank}_R S$, and $e>0$.
Thus $e=\operatorname{rank}_R S$ is equivalent to $f=1$ by cancellation.
An extension of fields has degree one precisely when its structure
map is bijective, equivalently surjective. The algebra $R[\pi]$
therefore maps onto $k_S$ and contains a uniformizer. The residue
and uniformizer generation lemma makes it all of $S$.
For the source-valued specialization, the chosen-extension theorem
proves the actual ring is a DVR and finite over the compatible-norm
unit ball. The identity of that ball with the original $R$ transports
finiteness and the actual inclusion maps. Its fraction field is $E$,
so the degree identity identifies its rank with $[E:K]$, and the
preceding argument applies directly over $R$.
\end{proof}

\begin{lemma}[An integral generator]
\label{fw:integral-generator}
\dcref{Serre1979, Chapter III, Section 6, Proposition 12}
\source{serre-1979-local-fields:page-0064}
\uses{fw:source-chosen-extension-generation,fw:totally-ramified-uniformizer-generation}
\lean{FontaineWintenberger.discreteValued_extension_exists_generator,FontaineWintenberger.discreteValued_extension_adjoin_uniformizer_eq_top}
\leanok
For every finite separable extension $E/K$ there is $x\in\mathcal O_E$ with
$\mathcal O_E=\mathcal O_K[x]$. If $E/K$ is totally ramified, any
uniformizer of $E$ is such a generator.
\end{lemma}
\begin{proof}
Choose a primitive element $\bar a$ of the finite separable extension
$k_E/k_K$. Lift its monic minimal polynomial to
$h\in\mathcal O_K[T]$. Hensel's lemma in $\mathcal O_E$ supplies a root
$a$ of $h$ reducing to $\bar a$, with $h'(a)$ a unit. For a uniformizer
$\pi$ of $E$, put $x=a+\pi$. Taylor expansion gives
$h(x)=h'(a)\pi+\pi^2b$ with $b\in\mathcal O_E$, so $z=h(x)$ is a
uniformizer and belongs to $A=\mathcal O_K[x]$.
The residue of $x$ generates $k_E$ over $k_K$.

Given $b_0\in\mathcal O_E$, choose $P_0\in\mathcal O_K[T]$ whose value
at $x$ has the residue of $b_0$. Write $b_0=P_0(x)+zb_1$ with
$b_1\in\mathcal O_E$ and repeat this procedure $e=e(E/K)$ times. It gives
\[
 b_0=\sum_{r=0}^{e-1}z^rP_r(x)+z^e b_e.
\]
Since $z^e\mathcal O_E=\mathfrak m_K\mathcal O_E$, we have
$\mathcal O_E=A+\mathfrak m_K\mathcal O_E$.
The quotient $\mathcal O_E/A$ is a finite $\mathcal O_K$-module, so
Nakayama's lemma makes it zero. If $E/K$ is totally ramified, the same
argument takes $A=\mathcal O_K[\pi]$ and $z=\pi$, since now
$k_E=k_K$. This proves the last assertion as well.
This is the residue-lifting proof of the cited proposition specialized
to complete fields, with Hensel's lemma used before the Taylor expansion
and Nakayama's lemma making the spanning step explicit.
\end{proof}

\begin{lemma}[Real lower groups from canonical integer groups]
\label{fw:real-lower-groups}
\dcref{Serre1979, Chapter IV, Sections 1 and 3}
\source{serre-1979-local-fields:page-0068,serre-1979-local-fields:page-0080}
\uses{fw:integer-lower-groups,fw:integral-order-function}
\lean{FontaineWintenberger.lowerRamificationGroupReal, FontaineWintenberger.lowerRamificationGroupReal_intCast, FontaineWintenberger.lowerRamificationGroupReal_antitone, FontaineWintenberger.lowerRamificationGroupReal_of_le_neg_one, FontaineWintenberger.lowerRamificationGroupReal_zero, FontaineWintenberger.lowerRamificationGroupReal_subgroupOf, FontaineWintenberger.lowerRamificationGroupReal_map_subtype, FontaineWintenberger.maximalIdeal_lowerRamificationGroupReal_normal, FontaineWintenberger.mem_lowerRamificationGroupReal_iff_le_order, FontaineWintenberger.lowerRamificationGroupReal_eventually_eq_bot}
\leanok
For a group acting on a commutative ring and an invariant ideal $I$,
put $D_t=D_{\lceil t\rceil}$ for every real $t$, using the canonical
integer congruence groups. This agrees with the integer groups at
integer indices, is decreasing, and commutes with restriction to a
subgroup. For a local ring with $I$ its maximal ideal, these groups are
normal. For a DVR with intrinsic order $j$, membership is equivalent to
$t+1\leq j(g)$ in the extended nonnegative reals. The groups are all of
$D$ when $t\leq-1$, and are eventually trivial for a finite faithful action.
\end{lemma}
\begin{proof}
The ceiling fixes integers and is increasing, giving agreement and
monotonicity from the integer results. Subgroup restriction and normality
follow at each integer ceiling. Integer membership is the bound
$\max(\lceil t\rceil+1,0)\leq j(g)$. If $j(g)=\infty$, both bounds hold.
Otherwise $j(g)=n$ is a nonnegative integer, and this is equivalent to
$\lceil t\rceil\leq n-1$, hence to $t+1\leq n$. Negative bounds are
automatic because $n\geq0$. If $t\leq-1$, then $\lceil t\rceil\leq-1$.
Finally an integer bound making all later integer groups trivial also
makes all real groups above that bound trivial, since $t\leq\lceil t\rceil$.
\end{proof}

\begin{lemma}[Finite upper ramification groups]
\label{fw:finite-upper-groups}
\dcref{Serre1979, Chapter IV, Section 3, Proposition 14}
\source{serre-1979-local-fields:page-0080,serre-1979-local-fields:page-0081}
\uses{fw:real-lower-groups,fw:herbrand-order-isomorphism}
\lean{FontaineWintenberger.upperRamificationGroup, FontaineWintenberger.upperRamificationGroup_normal, FontaineWintenberger.upperRamificationGroup_antitoneOn, FontaineWintenberger.upperRamificationGroup_neg_one, FontaineWintenberger.upperRamificationGroup_zero, FontaineWintenberger.upperRamificationGroup_herbrandFunction, FontaineWintenberger.upperRamificationGroup_eventually_eq_bot}
\leanok
For a finite group $G$ acting on a DVR, put
$G^u=G_{\psi_G(u)}$ for $u\geq-1$.
These are decreasing normal subgroups, with
$G^{-1}=G$, $G^0=G_0$, and $G^{\varphi_G(t)}=G_t$ for $t\geq-1$.
For a faithful action they are eventually trivial.
\end{lemma}
\begin{proof}
Normality follows from normality of the real lower groups. The inverse
Herbrand function is increasing, so reindexing a decreasing filtration
gives a decreasing filtration. The inverse is the identity on $[-1,0]$,
which proves the two specified values, and
$\psi_G(\varphi_G(t))=t$ proves the reindexing equality.
For a faithful action choose $T\geq-1$ above which the real lower groups
are trivial. If $u\geq\varphi_G(T)$, monotonicity gives
$\psi_G(u)\geq\psi_G(\varphi_G(T))=T$, proving eventual triviality.
\end{proof}

\begin{lemma}[Equivariant relabelling of ramification groups]
\label{fw:ramification-action-equiv}
\dcref{Serre1979, Chapter IV, Sections 1 and 3, definitions}
\source{serre-1979-local-fields:page-0068,serre-1979-local-fields:page-0080}
\uses{fw:integral-order-function,fw:herbrand-finite-sum,fw:finite-upper-groups}
\lean{FontaineWintenberger.ramificationOrder_mulEquiv, FontaineWintenberger.lowerRamificationGroup_map_mulEquiv, FontaineWintenberger.lowerRamificationGroupReal_map_mulEquiv, FontaineWintenberger.card_lowerRamificationGroupReal_mulEquiv, FontaineWintenberger.herbrandFunction_mulEquiv, FontaineWintenberger.inverseHerbrandFunction_mulEquiv, FontaineWintenberger.upperRamificationGroup_map_mulEquiv}
\leanok
Suppose $e:G\simeq G'$ identifies two actions on the same DVR, so
$e(g)(a)=g(a)$ for all $g,a$. Then $j_{G'}(e(g))=j_G(g)$, and $e$
maps every integer and real lower group onto the corresponding group.
If the groups are finite, their Herbrand functions and inverses agree
on $[-1,\infty)$, and $e(G^u)=G'^u$ for every $u\geq-1$.
\end{lemma}
\begin{proof}
Every integral displacement is unchanged, hence so is its infimum.
The order membership criterion identifies the lower subgroups, whose
cardinalities therefore agree. Their defining Herbrand integrals, including
the denominator at zero, agree. The inverse functions agree by uniqueness
of the inverse increasing bijection. Reindexing the identified real lower
groups by this common inverse proves the upper-group assertion.
\end{proof}

\begin{lemma}[Equivariant DVR isomorphisms preserve lower ramification]
\label{fw:ramification-dvr-lower-equiv}
\uses{fw:ideal-power-valuation,fw:integral-order-function,fw:real-lower-groups}
\lean{FontaineWintenberger.addVal_ringEquiv, FontaineWintenberger.ramificationOrder_equivariant, FontaineWintenberger.lowerRamificationGroupReal_map_equivariant}
\leanok
Let $r:A\simeq B$ be an isomorphism of DVRs and $e:G\simeq G'$ an
isomorphism of acting groups satisfying $r(g(a))=e(g)(r(a))$.
The normalized integral valuations satisfy $v_B(r(a))=v_A(a)$.
The intrinsic displacement orders agree under $e$, which maps the
actual real lower groups onto each other. No finiteness of the
acting groups is required.
\end{lemma}
\begin{proof}
The ring isomorphism maps the unique maximal ideal and each of its
powers onto its counterpart. The ideal-power valuation criterion
therefore gives the same lower bounds for the two normalized
valuations, including at zero, so those valuations agree.
Equivariance then identifies each integral displacement. Bijection
of the integral elements identifies their infima, proving the order
identity. The order membership criterion identifies the real lower
groups under $e$.
\end{proof}

\begin{lemma}[Equivariant DVR isomorphisms preserve ramification]
\label{fw:ramification-dvr-equiv}
\dcref{Serre1979, Chapter IV, Sections 1 and 3, definitions}
\source{serre-1979-local-fields:page-0068,serre-1979-local-fields:page-0080}
\uses{fw:ramification-dvr-lower-equiv,fw:finite-upper-groups}
\lean{FontaineWintenberger.card_lowerRamificationGroupReal_equivariant, FontaineWintenberger.herbrandFunction_equivariant, FontaineWintenberger.inverseHerbrandFunction_equivariant, FontaineWintenberger.upperRamificationGroup_map_equivariant}
\leanok
Let $r:A\simeq B$ be an isomorphism of DVRs and $e:G\simeq G'$ an
isomorphism of acting groups satisfying $r(g(a))=e(g)(r(a))$.
The normalized integral valuations satisfy $v_B(r(a))=v_A(a)$.
Consequently the intrinsic orders agree under $e$, which maps the real
lower groups onto each other. For finite groups the Herbrand functions,
inverse functions and upper groups agree under these identifications.
\end{lemma}
\begin{proof}
The ring isomorphism maps the unique maximal ideal, and hence every one
of its powers, onto its counterpart. The ideal-power criterion therefore
gives the same lower bounds for both normalized valuations, including
at zero, so the valuations agree. Equivariance identifies each integral
displacement, and the bijection of integral elements identifies their
infima. The lower membership criterion identifies the real lower groups.
Their cardinalities and the denominator at zero agree, so the defining
Herbrand integrals agree. Uniqueness identifies the inverses; reindexing
then identifies the upper groups.
\end{proof}

\begin{definition}[Ramification orders of general embeddings]
\label{fw:embedding-ramification-order}
\dcref{Wintenberger1983, 1.1.1}
\source{wintenberger}
\uses{fw:embedding-order-finite-ambient,fw:embedding-shifted-number}
\lean{FontaineWintenberger.closureEmbeddingOrder, FontaineWintenberger.closureEmbeddingOrder_eq_iInf, FontaineWintenberger.closureEmbeddingOrder_eq_top_iff, FontaineWintenberger.closureEmbeddingNumber, FontaineWintenberger.closureEmbeddingNumber_eq, FontaineWintenberger.closureEmbeddingNumber_add_one, FontaineWintenberger.closureEmbeddingNumber_self}
\leanok
For a finite separable extension $E/K$ and a $K$-embedding
$\sigma:E\longrightarrow K^{\mathrm{sep}}$, set
\[
 j_E(\sigma)=\inf_{a\in\mathcal O_E}v_E(\sigma(a)-a),
 \qquad i_E(\sigma)=j_E(\sigma)-1,
 \qquad j_E(\mathrm{id})=i_E(\mathrm{id})=+\infty.
\]
Here the normalized valuation is extended from $E$ to the separable
closure. The image of $\sigma$ need not lie in $E$.
\end{definition}

\begin{definition}[Orders of pairs in a finite ambient field]
\label{fw:finite-embedding-pair-order}
\dcref{Wintenberger1983, 1.1.1}
\source{wintenberger}
\uses{fw:normalized-field-valuation,fw:finite-valuation-ring,fw:spectral-integer-residue-maps,fw:ramification-index-scaling}
\lean{FontaineWintenberger.embeddingRamificationIdx, FontaineWintenberger.embeddingIntegralOrder, FontaineWintenberger.embeddingRamificationOrder, FontaineWintenberger.embeddingRamificationIdx_pos, FontaineWintenberger.embeddingIntegralOrder_eq_top_iff, FontaineWintenberger.embeddingRamificationOrder_eq_top_iff, FontaineWintenberger.embeddingRamificationOrder_self, FontaineWintenberger.embeddingRamificationOrder_lt_top}
\leanok
Let $K$ be a complete discretely valued field with perfect residue
field of characteristic $p>0$. Let $E,N$ be finite separable
extensions of $K$, and let $\sigma,\tau:E\longrightarrow N$ be
$K$-embeddings. Write $\sigma^\circ,\tau^\circ$ for their actual
maps on valuation rings and $e_\sigma=e(N/\sigma E)$ for the
ramification index of $\sigma^\circ$. Define
\[
 r_N(\sigma,\tau)=\inf_{a\in\mathcal O_E}
 v_N(\sigma^\circ(a)-\tau^\circ(a))\in\mathbb N\cup\{\infty\},
 \qquad
 J_{E,N}(\sigma,\tau)=\frac{r_N(\sigma,\tau)}{e_\sigma}
 \in[0,\infty].
\]
Here $v_N$ is the normalized additive integral valuation, and
$\infty/e_\sigma=\infty$; the denominator is a positive integer.
The definition uses the given finite fields and their original
valuation rings. It requires neither field to be normal over $K$.
\end{definition}

\begin{lemma}[Evaluation of an embedding pair on an integral generator]
\label{fw:embedding-pair-generator}
\dcref{Serre1979, Chapter IV, Section 1, Lemma 1}
\source{serre-1979-local-fields:page-0068,serre-1979-local-fields:page-0069}
\uses{fw:finite-embedding-pair-order,fw:henselian-integral-generation,fw:ideal-power-valuation}
\lean{FontaineWintenberger.exists_embeddingIntegralOrder_generator, FontaineWintenberger.embeddingIntegralOrder_eq_generator, FontaineWintenberger.embeddingRamificationOrder_eq_generator}
\leanok
There exists $x\in\mathcal O_E$ with
$\mathcal O_E=\mathcal O_K[x]$. Every such $x$ satisfies
\[
 r_N(\sigma,\tau)=v_N(\sigma(x)-\tau(x)),\qquad
 J_{E,N}(\sigma,\tau)=\frac{v_N(\sigma(x)-\tau(x))}{e_\sigma}.
\]
In particular $J_{E,N}(\sigma,\tau)=\infty$ if and only if
$\sigma=\tau$; distinct embeddings have finite order.
\end{lemma}
\begin{proof}
The integral-generation theorem supplies $x$: the base is complete
and hence henselian, its residue field is perfect, and $E/K$ is
finite separable. Fix any integral generator $x$ and an ideal
$I\subseteq\mathcal O_N$. If
$\sigma(x)-\tau(x)\in I$, the two maps to $\mathcal O_N/I$ agree
on $x$ and on $\mathcal O_K$, and therefore on every element of
$\mathcal O_K[x]$. Explicitly, the property that their difference
belongs to $I$ contains the base and $x$, is closed under addition,
and is closed under multiplication by
\[
 \sigma(ab)-\tau(ab)
 =\sigma(a)(\sigma(b)-\tau(b))
   +(\sigma(a)-\tau(a))\tau(b).
\]
The converse follows by evaluating at $x$. Apply this equivalence
to $I=\mathfrak m_N^m$ for each $m\geq0$. The ideal-power
valuation criterion says that $r_N(\sigma,\tau)$ and
$v_N(\sigma(x)-\tau(x))$ have precisely the same natural-number
lower bounds, so they are equal, also at infinity. Division gives
the formula for $J$.

If $\sigma=\tau$, every displacement is zero and has value
$\infty$. Conversely an infinite order makes
$\sigma(x)=\tau(x)$, hence the two maps agree on all of
$\mathcal O_E$ by the generator argument with $I=0$.
Every element of $E$ is a quotient of two elements of
$\mathcal O_E$ with nonzero denominator; the field embeddings
therefore agree on $E$. If they are distinct, the generator
displacement is nonzero, and its integral valuation is a finite
natural number. The index is positive and finite, so the normalized
order is finite as well.
\end{proof}

\begin{lemma}[Independence of the finite ambient field]
\label{fw:embedding-pair-ambient-invariance}
\dcref{Wintenberger1983, 1.1.1, normalization of the extended valuation}
\source{wintenberger}
\uses{fw:embedding-pair-generator,fw:normalized-valuation-scaling,fw:spectral-integer-residue-maps}
\lean{FontaineWintenberger.embeddingRamificationIdx_comp, FontaineWintenberger.embeddingRamificationIdx_eq, FontaineWintenberger.embeddingIntegralOrder_comp, FontaineWintenberger.embeddingRamificationOrder_comp, FontaineWintenberger.embeddingIntegralOrder_comm, FontaineWintenberger.embeddingRamificationOrder_comm}
\leanok
Let $M/K$ be finite separable and $u:N\longrightarrow M$ a
$K$-embedding. Then
\[
 r_M(u\sigma,u\tau)=e(M/uN)r_N(\sigma,\tau),\qquad
 J_{E,M}(u\sigma,u\tau)=J_{E,N}(\sigma,\tau).
\]
Also $e(N/\sigma E)=e(N/\tau E)$, so the normalized pair order
is symmetric in $\sigma,\tau$.
\end{lemma}
\begin{proof}
Choose an actual integral generator $x$ of $\mathcal O_E$.
The integer map induced by $u$ is injective and local. The
normalized valuation scaling theorem gives
\[
 v_M(u(\sigma(x)-\tau(x)))
 =e(M/uN)v_N(\sigma(x)-\tau(x)).
\]
The generator formula proves the raw-order identity. The actual
tower of integer maps from $\mathcal O_E$ through
$\mathcal O_N$ to $\mathcal O_M$ has index
\[
 e(M/u\sigma E)=e(N/\sigma E)e(M/uN).
\]
These are positive finite integers. Cancelling $e(M/uN)$ in the
two displayed formulas proves normalized invariance; if the raw
order is infinite, both normalized orders are infinite, and the
same conclusion holds. Finally the two actual towers from
$\mathcal O_K$ through $\mathcal O_E$ to $\mathcal O_N$ give
\[
 e(N/K)=e(E/K)e(N/\sigma E)
       =e(E/K)e(N/\tau E).
\]
Since $e(E/K)>0$, cancellation identifies the two relative indices.
The valuation of a negative equals that of the original element,
so interchanging $\sigma,\tau$ preserves the raw order and now
also the normalized order.
\end{proof}

\begin{lemma}[The closure order is the finite-pair order]
\label{fw:embedding-order-finite-ambient}
\dcref{Wintenberger1983, 1.1.1}
\source{wintenberger}
\uses{fw:embedding-pair-ambient-invariance,fw:normalized-valuation-scaling}
\lean{FontaineWintenberger.embeddingPairField, FontaineWintenberger.embeddingPairLeft, FontaineWintenberger.embeddingPairRight, FontaineWintenberger.closureEmbeddingOrder, FontaineWintenberger.closureEmbeddingOrder_eq_iInf, FontaineWintenberger.closureEmbeddingOrder_eq_finite, FontaineWintenberger.closureEmbeddingOrder_comp, FontaineWintenberger.closureEmbeddingOrder_comm}
\leanok
For two $K$-embeddings $\sigma,\tau:E\longrightarrow K^{\rm sep}$,
their images generate a finite separable field
$N=\sigma(E)\tau(E)$. Compute their pair order in $N$ after
factoring both maps through $N$. This value is unchanged if $N$
is replaced by any finite separable subfield containing both
images. For comparison with a fixed inclusion
$\eta:E\longrightarrow K^{\rm sep}$, take $\tau=\eta$ and
$N=\eta(E)\sigma(E)$. Then
\[
 j_E(\sigma)=J_{E,N}(\eta,\sigma).
\]
Thus this finite-ambient definition realizes the original
closure-valued displacement infimum. It is infinite precisely
when the two closure embeddings agree.
\end{lemma}
\begin{proof}
Both images have finite degree over $K$, and their compositum has
degree at most the product of those degrees. It is separable
because it lies in $K^{\rm sep}$. If $N_1,N_2$ are two finite
ambient subfields containing the images, their compositum is
again finite separable; its inclusion maps and the preceding
invariance theorem identify both computations with the same
one in $N_1N_2$.

The normalized valuation of $N$, divided by $e(N/\eta E)$,
restricts to the normalized valuation of $E$, by valuation
scaling. These divided valuations agree whenever a finite
ambient field is enlarged, again by scaling and index
multiplicativity. They therefore define the extension of that
valuation on the union of finite subfields containing $\eta(E)$, namely
$K^{\rm sep}$. For every $a\in\mathcal O_E$ the resulting
value of $\sigma(a)-\eta(a)$ is exactly its value in $N$
divided by $e(N/\eta E)$. The generator formula computes the
infimum at the same actual integral generator in both
descriptions, proving the equality of orders. The infinity
criterion follows from the finite-pair criterion and injectivity
of the inclusion of $N$ into the closure.
\end{proof}

\begin{lemma}[Residue and uniformizer criteria for embedding pairs]
\label{fw:embedding-pair-uniformizer}
\dcref{Wintenberger1983, 1.1.1; Serre1979, Chapter IV, Section 2, Proposition 5}
\source{wintenberger,serre-1979-local-fields:page-0073}
\uses{fw:finite-embedding-pair-order,fw:finite-perfect-residue-digits,fw:spectral-integer-residue-maps,fw:ideal-power-valuation}
\lean{FontaineWintenberger.one_le_embeddingIntegralOrder_iff, FontaineWintenberger.embeddingIntegralOrder_eq_zero_of_residue_ne, FontaineWintenberger.embeddingRamificationOrder_eq_zero_of_residue_ne, FontaineWintenberger.embeddingIntegralOrder_eq_uniformizer, FontaineWintenberger.embeddingRamificationOrder_eq_uniformizer}
\leanok
The induced residue embeddings $\bar\sigma,\bar\tau:k_E\to k_N$
agree if and only if $r_N(\sigma,\tau)\geq1$. If they differ,
$r_N(\sigma,\tau)=J_{E,N}(\sigma,\tau)=0$. If they agree, then
for every uniformizer $\pi$ of $E$,
\[
 r_N(\sigma,\tau)=v_N(\sigma(\pi)-\tau(\pi)),\qquad
 J_{E,N}(\sigma,\tau)
 =\frac{v_N(\sigma(\pi)-\tau(\pi))}{e(N/\sigma E)}.
\]
These formulas include equal embeddings and residue characteristic
$2$, and require neither a finite residue field nor normality.
\end{lemma}
\begin{proof}
Reduction of each integral displacement is
$\bar\sigma(\bar a)-\bar\tau(\bar a)$. Surjectivity of
$\mathcal O_E\to k_E$ and the maximal-ideal valuation
criterion prove the first equivalence. Its failure gives a unit
displacement, so the infimum is zero, as is its quotient by the
positive index.

Suppose the residue maps agree and fix $m\geq1$. Write
$f=\sigma^\circ$, $g=\tau^\circ$, and
$I=\mathfrak m_N$. If $f(\pi)-g(\pi)\in I^m$, we show
$f(a)-g(a)\in I^m$ for every $a\in\mathcal O_E$.
First, whenever $z-w\in I^r$ with $r\geq1$ and $p\in I$,
the binomial expansion gives $z^p-w^p\in I^{r+1}$:
the term linear in $z-w$ contains $p$, and every term of
degree at least two belongs to $I^{2r}\subseteq I^{r+1}$.
Induction therefore gives
\[
 z-w\in I\quad\Longrightarrow\quad
 z^{p^n}-w^{p^n}\in I^{n+1}\quad(n\geq0).
\]
Use the finite perfect-residue digit theorem with $n=r=m$:
\[
 a=\sum_{j<m}b_j^{p^m}\pi^j+\pi^m c.
\]
For each $b_j$, equality of residue maps puts
$f(b_j)-g(b_j)$ in $I$, so the displayed power congruence
puts their $p^m$-power difference in
$I^{m+1}\subseteq I^m$. The identity for a difference of
$j$th powers makes $f(\pi)^j-g(\pi)^j$ a multiple of
$f(\pi)-g(\pi)$, hence an element of $I^m$. The product
difference identity used in the generator proof now puts the
difference of each digit term in $I^m$. Localness puts both
$f(\pi)$ and $g(\pi)$ in $I$, so the images of the two
remainders also belong to $I^m$. Adding proves the assertion
for $a$. For $m=0$ it holds because $I^0=\mathcal O_N$.
The converse follows by evaluating at $\pi$. Thus the raw
infimum and the uniformizer displacement have the same
natural-number lower bounds, proving their equality, even at
infinity. Division proves the normalized formula. This adapts
the cited uniformizer criterion to two embeddings by finite
perfect-residue digits rather than a Galois automorphism argument.
\end{proof}

\begin{lemma}[Agreement with the original automorphism order]
\label{fw:embedding-pair-automorphism}
\dcref{Wintenberger1983, 1.1.1, final remark}
\source{wintenberger}
\uses{fw:embedding-pair-ambient-invariance,fw:integral-order-function,fw:spectral-full-galois-action}
\lean{FontaineWintenberger.embeddingIntegralOrder_aut, FontaineWintenberger.embeddingRamificationOrder_aut, FontaineWintenberger.closureEmbeddingOrder_aut}
\leanok
For a $K$-automorphism $s$ of a finite separable extension $E/K$,
\[
 r_E(s,\mathrm{id})=j_E(s),\qquad
 J_{E,E}(s,\mathrm{id})=j_E(s),
\]
where the right-hand side is the original intrinsic automorphism
ramification order, viewed in $[0,\infty]$ for the second equality.
No normality assumption on $E/K$ is imposed.
\end{lemma}
\begin{proof}
The actual spectral integer map of $s$ is exactly its original
action on $\mathcal O_E$. Both raw-order definitions are therefore
the infimum of the same integral displacements. Independence of the
relative index from the source embedding identifies the index of
$s$ with that of the identity. Composing the identity with itself
and multiplying indices gives $e=e^2$ for this positive integer,
so $e=1$. Division by this index leaves the raw order unchanged,
proving the normalized equality.
\end{proof}

\begin{lemma}[Signed embedding ramification numbers]
\label{fw:embedding-shifted-number}
\dcref{Wintenberger1983, 1.1.1}
\source{wintenberger}
\uses{fw:embedding-order-finite-ambient,fw:embedding-pair-uniformizer,fw:embedding-pair-automorphism}
\lean{FontaineWintenberger.embeddingRamificationNumber, FontaineWintenberger.embeddingRamificationNumber_add_one, FontaineWintenberger.neg_one_le_embeddingRamificationNumber, FontaineWintenberger.embeddingRamificationNumber_ne_bot, FontaineWintenberger.embeddingRamificationNumber_eq_top_iff, FontaineWintenberger.embeddingRamificationNumber_exists_real, FontaineWintenberger.embeddingRamificationNumber_eq_neg_one_of_residue_ne, FontaineWintenberger.embeddingRamificationNumber_comp, FontaineWintenberger.embeddingRamificationNumber_comm, FontaineWintenberger.embeddingRamificationNumber_nonneg, FontaineWintenberger.embeddingRamificationNumber_eq_uniformizer, FontaineWintenberger.closureEmbeddingNumber_comp, FontaineWintenberger.closureEmbeddingNumber_aut}
\leanok
For a pair of finite-field embeddings define
\[
 I_{E,N}(\sigma,\tau)=J_{E,N}(\sigma,\tau)-1
 \quad\hbox{in }\mathbb R\cup\{-\infty,+\infty\},
\]
using the canonical inclusion of $[0,\infty]$ into the extended
real line. Define the closure-pair number in the same way from
its finite-ambient order. Both numbers belong to $[-1,+\infty]$,
and adding $1$ recovers the original nonnegative order exactly.
They are infinite precisely for equal embeddings; distinct
embeddings have finite real numbers. They are symmetric and
unchanged by finite ambient enlargement. If the residue embeddings
differ the number is $-1$; if they agree it is nonnegative and is the normalized
uniformizer-displacement value minus $1$, for every source
uniformizer. For an automorphism it is the original unshifted
automorphism order minus $1$. In particular these numbers give
the signed $i_E$ and the identity endpoint in the original
closure-embedding definition.
\end{lemma}
\begin{proof}
The canonical extended-real inclusion preserves order and sends
$0$ to $0$ and $+\infty$ to $+\infty$, never to $-\infty$.
Subtracting the finite number $1$ therefore gives a value at
least $-1$; it preserves finiteness and the positive-infinity
endpoint, and addition of $1$ cancels that subtraction even at
$+\infty$. The proved infinity criterion for $J$ identifies
equal embeddings and shows that the remaining values are finite
real numbers. Symmetry and ambient independence follow by
subtracting $1$ from the corresponding equalities for $J$.
Unequal residue maps have $J=0$ and hence $I=-1$. Equal residue
maps give the proved uniformizer formula for $J$. Both images of
a source uniformizer have target valuation $e(N/\sigma E)$,
by valuation scaling and independence of this index from the
embedding. The valuation inequality makes their difference have
at least that value. Thus $J\geq1$ and $I\geq0$; subtraction
also gives the claimed uniformizer formula for $I$. Applying the actual
automorphism comparison to $J$ proves the last assertion. All
statements for closure pairs follow from their actual finite
image compositum and the proved comparisons with every finite
ambient field containing their images.
\end{proof}

\begin{definition}[Lower and upper ramification numbering]
\label{fw:ramification-filtration}
\dcref{Serre1979, Chapter IV, Sections 1 and 3; Wintenberger1983, 1.1.1--1.1.2}
\source{serre-1979-local-fields:page-0068,serre-1979-local-fields:page-0080}
\uses{fw:integer-lower-groups,fw:integral-order-function,fw:real-lower-groups,fw:finite-upper-groups,fw:infinite-upper-groups}
\lean{FontaineWintenberger.ramificationOrder, FontaineWintenberger.ramificationOrder_one, FontaineWintenberger.mem_lowerRamificationGroup_iff_le_order, FontaineWintenberger.lowerRamificationGroupReal, FontaineWintenberger.herbrandFunction, FontaineWintenberger.inverseHerbrandFunction, FontaineWintenberger.upperRamificationGroup, FontaineWintenberger.absoluteUpperRamificationGroup}
\leanok
For a finite Galois extension $E/K$ with group $G$, put
$j_E(s)=\inf_{a\in\mathcal O_E}v_E(s(a)-a)$ and $j_E(1)=+\infty$.
For finite Galois $E/K$ with group $G$, its lower groups are
\[
 G_{-1}=G,\qquad
 G_n=\{s\in G:j_E(s)\geq n+1\}\quad(n\in\mathbb Z_{\geq0}),
 \qquad G_t=G_{\lceil t\rceil}\quad(t\geq-1).
\]
Define
\[
 \varphi_{E/K}(t)=
 \begin{cases}
 t,&-1\leq t\leq0,\\
 \displaystyle\int_0^t\frac{|G_s|}{|G_0|}\,ds,&t\geq0,
 \end{cases}
 \qquad \psi_{E/K}=\varphi_{E/K}^{-1},\qquad
 G^u=G_{\psi_{E/K}(u)}.
\]
Existence of the inverse is proved below. The absolute upper groups are
\[
 G_K^u=\bigcap_{E/K\ \mathrm{finite\ Galois}}
    \operatorname{res}_{E/K}^{-1}(\Gal(E/K)^u).
\]
\end{definition}

\begin{lemma}[Congruences detected on an algebra generator]
\label{fw:generator-evaluation}
\dcref{Serre1979, Chapter IV, Section 1, Lemma 1}
\source{serre-1979-local-fields:page-0068,serre-1979-local-fields:page-0069}
\uses{fw:henselian-integral-generation,fw:integral-order-function,fw:spectral-full-galois-action}
\lean{FontaineWintenberger.mem_inertia_iff_generator, FontaineWintenberger.ramificationOrder_eq_generator, FontaineWintenberger.mem_lowerRamificationGroup_iff_generator, FontaineWintenberger.spectralIntegerRing_decompositionSubgroup_smulCommClass, FontaineWintenberger.spectralIntegerRing_exists_order_generator, FontaineWintenberger.spectralIntegerRing_exists_galois_order_generator}
\leanok
Let a group $D$ act on a commutative $R$-algebra $S$ by algebra
automorphisms, and suppose $S=R[x]$. For every ideal $I$ of $S$,
\[
 g\in I.\operatorname{inertia}(D)
 \quad\Longleftrightarrow\quad g(x)-x\in I.
\]
If $S$ is a discrete valuation ring, its intrinsic order function is
$j_S(g)=v_S(g(x)-x)$. In particular every integer lower congruence group
is detected by the corresponding bound at $x$.
In the spectral extension setting of the Henselian integral-generation
lemma, such an element $x$ exists simultaneously for every element of
the intrinsic decomposition group and, under the preceding full-group
identification, for every $K$-automorphism of $E$.
\end{lemma}
\begin{proof}
The forward implication evaluates the defining congruence at $x$.
For the converse, the elements $y$ with $g(y)-y\in I$ contain the
image of $R$ and $x$, and are closed under addition and multiplication:
\[
 g(yz)-yz=(g(y)-y)g(z)+y(g(z)-z).
\]
They therefore contain $R[x]=S$. Apply this equivalence to every power
of the maximal ideal. The valuation/ideal-power criterion says precisely
that every natural lower bound for $v_S(g(x)-x)$ is a lower bound for
all integral displacements. Equality with their infimum follows, also
when the value is infinite.
For the spectral assertion, choose $x$ from the Henselian
integral-generation lemma. Decomposition automorphisms fix base integers,
so the equality applies to every such automorphism with this same $x$.
\end{proof}

\begin{lemma}[Order function and lower filtration]
\label{fw:lower-ramification-properties}
\dcref{Serre1979, Chapter IV, Section 1, Lemma 1 and Proposition 1}
\source{serre-1979-local-fields:page-0068,serre-1979-local-fields:page-0069}
\uses{fw:integral-generator,fw:ramification-filtration,fw:integral-order-function}
\lean{FontaineWintenberger.ramificationOrder_eq_generator, FontaineWintenberger.ramificationOrder_eq_top_iff, FontaineWintenberger.lowerRamificationGroupReal_antitone, FontaineWintenberger.maximalIdeal_lowerRamificationGroupReal_normal, FontaineWintenberger.lowerRamificationGroupReal_eventually_eq_bot, FontaineWintenberger.ramificationOrder_inv, FontaineWintenberger.ramificationOrder_conj, FontaineWintenberger.min_order_le_order_mul, FontaineWintenberger.ramificationOrder_mul_of_ne}
\leanok
For finite Galois $E/K$ and $\mathcal O_E=\mathcal O_K[x]$,
$j_E(s)=v_E(s(x)-x)$. In particular $j_E(s)$ is a nonnegative integer
for $s\ne1$. The $G_t$ form a decreasing family of normal subgroups,
$G_t=\{1\}$ for all sufficiently large $t$, and
\[
 j_E(s^{-1})=j_E(s),\qquad j_E(rsr^{-1})=j_E(s),\qquad
 j_E(st)\geq\min\{j_E(s),j_E(t)\}.
\]
The last inequality is an equality if $j_E(s)\ne j_E(t)$.
\end{lemma}
\begin{proof}
For any polynomial $P\in\mathcal O_K[T]$, the difference
$P(s(x))-P(x)$ is divisible by $s(x)-x$ in $\mathcal O_E$.
This proves the formula for $j_E$. An automorphism fixing $x$ fixes $E$,
so its value is finite unless it is the identity.
Every $K$-automorphism preserves the valuation, by uniqueness of its
extension from $K$. For $n\geq0$, $G_n$ is the kernel of the action
on $\mathcal O_E/\mathfrak m_E^{n+1}$; hence it is normal.
The groups decrease, and choosing an integer larger than all the finitely
many $j_E(s)$ for $s\ne1$ proves eventual triviality.

Replacing $a$ by $s^{-1}(a)$ in the infimum proves the inverse identity;
replacing it by $r^{-1}(a)$ proves conjugacy invariance. The equality
$st(a)-a=s(t(a)-a)+(s(a)-a)$ gives the product inequality.
If $j_E(s)<j_E(t)$ and $j_E(st)>j_E(s)$, applying that inequality to
$s=(st)t^{-1}$ is a contradiction. The case $j_E(t)<j_E(s)$ follows by
applying it to $t=s^{-1}(st)$.
\end{proof}

\begin{lemma}[Lower numbering for subgroups]
\label{fw:ramification-subgroups}
\dcref{Serre1979, Chapter IV, Section 1, Proposition 2}
\source{serre-1979-local-fields:page-0069}
\uses{fw:spectral-base-integers,fw:ramification-dvr-lower-equiv,fw:integer-lower-groups,fw:real-lower-groups,fw:integral-order-function}
\lean{FontaineWintenberger.ramificationOrder_subgroup, FontaineWintenberger.lowerRamificationGroup_subgroupOf, FontaineWintenberger.lowerRamificationGroup_map_subtype, FontaineWintenberger.lowerRamificationGroupReal_subgroupOf, FontaineWintenberger.lowerRamificationGroupReal_map_subtype, FontaineWintenberger.finiteBase_ramificationOrder, FontaineWintenberger.finiteBase_lowerReal_map_fixingSubgroupEquiv, FontaineWintenberger.finiteBase_lower_map_fixingSubgroupEquiv, FontaineWintenberger.finiteBase_lowerReal_map_restrictScalars, FontaineWintenberger.finiteBase_lower_map_restrictScalars}
\leanok
If $H\subset G=\Gal(E/K)$ and $F=E^H$, the order function for $E/F$
is the restriction of $j_E$ and $H_t=H\cap G_t$ for every $t\geq-1$.
\end{lemma}
\begin{proof}
Give $E$ its canonical extended norm and $F$ the restricted norm.
The spectral integer-ring comparison identifies the actual integer rings
constructed over $K$ and over $F$ by the identity on $E$. Its
equivariance identifies the actual fixing-subgroup action with the
$F$-automorphism action. The equivariant DVR isomorphism preserves the
normalized integral valuation. Thus both order functions use the same
valuation on $E$ and the same set $\mathcal O_E$. Their infima agree for each $s\in H$.
The integer-indexed groups therefore agree, and so do their extensions
by the ceiling convention.
\end{proof}

\begin{theorem}[Galois descent to integral rings and residue inertia]
\label{fw:galois-ring-residue-theory}
\dcref{MathlibGaloisInertia, Galois descent, residue stabilizers and inertia cardinality}
\source{mathlib-is-galois-group,mathlib-galois-group-degree,mathlib-invariant-basic,mathlib-invariant-galois,mathlib-galois-ramification}
\lean{IsGaloisGroup.of_isFractionRing,IsGaloisGroup.card_eq_finrank',Ideal.Quotient.stabilizerHom_surjective,Ideal.Quotient.stabilizerQuotientInertiaEquiv,Ideal.Quotient.normal,Ideal.card_inertia_eq_ramificationIdxIn,Ideal.ramificationIdxIn_eq_ramificationIdx}
\mathlibok
Let a finite group $G$ act faithfully on a domain $S$, fixing exactly
the image of a domain $R$. If $R\hookrightarrow S$ is injective, then
$|G|=\operatorname{rank}_R S$. For maximal ideals $\mathfrak q$ over $\mathfrak p$,
the residue extension is normal, and its automorphism group is the
quotient of the stabilizer of $\mathfrak q$ by ideal inertia.
If $S$ is finite flat over $R$ and the residue extensions over
$\mathfrak p$ are separable, then
$|I_{\mathfrak q}|=e(\mathfrak q/\mathfrak p)$ for the canonical
local-length ramification index.
For compatible fraction fields $E/K$, a faithful Galois action on
$E/K$ preserving $S$ descends to such a Galois action on $S/R$
when $S$ is integral over the integrally closed ring $R$.
\end{theorem}
\begin{proof}
These are the invariant-ring theorems of \cite{MathlibGaloisInertia}.
Orbit polynomials have invariant coefficients, hence coefficients in
$R$; their reductions give normality. The stabilizer reduction theorem
makes the residue action surjective, and its kernel consists precisely
of the automorphisms whose differences lie in $\mathfrak q$, namely
ideal inertia. Thus the first isomorphism theorem gives the stated
quotient. With separable residue fields its order is the residue
degree. Transitivity on primes above $\mathfrak p$ makes all their
ramification indices and residue degrees equal. Orbit-stabilizer and
the finite-flat degree identity then give, respectively,
$|G|=g|I_{\mathfrak q}|f$ and $|G|=gef$.
Here $g$ and $f$ are positive, so cancellation proves the inertia formula.
The degree assertion is the fixed-field degree theorem after passing to
fraction fields; localization preserves module rank.
For descent, faithfulness on $S$ follows because fractions of $S$
generate $E$. A fixed element of $S$ belongs to $K$ and is integral
over $R$; integral closedness puts it in $R$. This proves the fixed-ring
assertion and hence descent. The source proves the residue stabilizer
surjectivity and normality in this invariant-ring generality; its
separability requirement enters only the conversion of automorphism
order to residue degree.
\end{proof}

\begin{lemma}[Inertia size on actual finite-extension integers]
\label{fw:local-inertia-size}
\dcref{Serre1979, Chapter IV, Section 1, Proposition 1}
\source{serre-1979-local-fields:page-0069}
\uses{fw:galois-ring-residue-theory,fw:spectral-full-galois-action,fw:finite-valuation-ring,fw:integer-lower-groups}
\lean{FontaineWintenberger.hasSeparableResidueFieldsAt_maximalIdeal, FontaineWintenberger.card_lowerRamificationGroup_zero, FontaineWintenberger.localResidueGaloisHom, FontaineWintenberger.localResidueGaloisHom_residue, FontaineWintenberger.localResidueGaloisHom_surjective, FontaineWintenberger.ker_localResidueGaloisHom, FontaineWintenberger.localResidueGaloisEquiv, FontaineWintenberger.localResidueGaloisEquiv_mk, FontaineWintenberger.localResidue_normal, FontaineWintenberger.spectralIntegerRing_galois_smulDistribClass, FontaineWintenberger.spectralIntegerRing_isGaloisGroup, FontaineWintenberger.spectralIntegerRing_residue_isGalois, FontaineWintenberger.spectralIntegerRing_card_lower_zero}
\leanok
For a finite injective local map of DVRs $R\hookrightarrow S$ with
faithful finite group $G$ fixing exactly $R$, reduction gives a
surjective homomorphism $G\to\operatorname{Aut}(k_S/k_R)$ whose kernel
is the zero lower group; it induces
$G/G_0\simeq\operatorname{Aut}(k_S/k_R)$, and the residue extension is normal.
If the residue extensions are separable, the zero lower group has order
$e(S/R)$.
In particular, if $E/K$ is finite Galois and $K$ is a complete
nonarchimedean normed field with DVR unit ball and perfect residue
field, the actual spectral integer ring, under its full field
automorphism action, has Galois residue extension and zero lower group
of order $e(E/K)$.
\end{lemma}
\begin{proof}
Every automorphism of the local ring $S$ preserves its maximal ideal,
so its stabilizer is all of $G$. The preceding residue theorem makes
reduction onto the residue automorphism group surjective, with ideal
inertia as kernel, and gives residue normality. The zero lower group
is this ideal inertia. For its order, apply the preceding inertia theorem to this ideal, lying over the
maximal ideal of $R$ by localness, and use its canonical index.
For the specialization, the actual integer ring is integral and finite
free over the base unit ball, hence flat. Its full automorphism action
agrees with the field action on elements. Galois descent therefore
proves that its fixed ring is exactly the base unit ball. Perfectness
of the base residue field makes every residue extension separable,
so the general assertion applies.
\end{proof}

\begin{lemma}[The inertia-relative index of a nonnormal subextension]
\label{fw:inertia-relative-index}
\uses{fw:local-inertia-size,fw:fixed-field-integers,fw:coset-ramification-average,fw:normalized-valuation-scaling}
\lean{FontaineWintenberger.fixedField_ramificationIdx_eq_relIndex, FontaineWintenberger.fixedField_sup_inertia_eq_top, FontaineWintenberger.fixedField_ramificationIdx_eq_finrank_iff}
\leanok
Let $\Omega/K$ be finite Galois over a complete discretely valued
field with perfect residue field. Write $G=\Gal(\Omega/K)$,
$D=G_0$, and let $H\subseteq G$ be any subgroup, with actual
fixed field $F=\Omega^H$. Then
\[
 e(F/K)=[D:H\cap D].
\]
In particular $F/K$ is totally ramified exactly when $HD=G$.
Every valuation ring in these formulas is the actual restriction
of the ambient valuation; it also equals the actual spectral ring.
\end{lemma}
\begin{proof}
Let $q=e(\Omega/F)$. The actual restricted $H$-action has inertia
$H\cap D$, so $q=|H\cap D|$. The full action has
$e(\Omega/K)=|D|$. Multiplicativity of ramification indices gives
$e(F/K)q=|D|$. Lagrange's formula gives
$[D:H\cap D]q=|D|$, and $q>0$ permits cancellation.
Since $D$ is normal, multiplication defines the subgroup $HD$.
The inclusion of $D$ identifies its cosets modulo $H\cap D$
with the cosets of $H$ in $HD$: surjectivity follows by writing
each element as $hd$, and equality of two cosets is exactly
membership of the quotient in $H\cap D$. Thus
\[
 [G:H]=[G:HD][D:H\cap D].
\]
The actual fixed-field degree is $[F:K]=[G:H]$. Thus
$e(F/K)=[F:K]$ holds exactly when $[G:HD]=1$, namely $HD=G$.
Equality of the actual restricted
and spectral valuation rings follows from uniqueness of the
extension valuation ring.
\end{proof}

\begin{lemma}[Total ramification in a local tower]
\label{fw:total-ramification-tower}
\uses{fw:local-dvr-degree-identity}
\lean{FontaineWintenberger.ramificationIdx_eq_finrank_tower_iff}
\leanok
Let $R\subseteq S\subseteq T$ be compatible injective local maps
of DVRs, finite as modules, and write $n(S/R)$ for their module
ranks. Then
\[
 e(T/R)=n(T/R)\quad\Longleftrightarrow\quad
 e(S/R)=n(S/R)\text{ and }e(T/S)=n(T/S).
\]
For their fraction fields these are the ordinary field degrees.
In particular every step of an actual totally ramified finite
extension is totally ramified.
\end{lemma}
\begin{proof}
The local degree identity identifies equality of ramification index
and rank with surjectivity of the canonical residue-field map.
The residue map $k_R\to k_T$ is the composite of the maps
$k_R\to k_S\to k_T$, which are field embeddings. If the composite
is surjective, the second map is surjective. For $s\in k_S$,
choose $r\in k_R$ with the same image in $k_T$; injectivity of
$k_S\to k_T$ then gives equality of its image with $s$. Thus the
first map is surjective too. Conversely two surjective maps have
surjective composite. Finally localization identifies each module
rank with the corresponding fraction-field degree.
\end{proof}

\begin{corollary}[The residue index of a nonnormal subextension]
\label{fw:inertia-residue-index}
\uses{fw:local-inertia-size,fw:ramification-subgroups,fw:normalized-valuation-scaling,fw:local-dvr-degree-identity}
\lean{FontaineWintenberger.intermediate_spectral_zeroOrbitIndex_eq_residueDegree}
\leanok
In the preceding setting, $f(F/K)=[G:HD]$.
\end{corollary}
\begin{proof}
The coset formula in the preceding proof gives
$[F:K]=[G:HD]e(F/K)$. The fundamental degree identity gives
$[F:K]=f(F/K)e(F/K)$. Cancellation of the positive integer
$e(F/K)$ proves the equality.
\end{proof}

\begin{lemma}[The original field fixed by inertia]
\label{fw:inertia-fixed-field}
\source{serre-1979-local-fields:page-0069}
\dcref{A consequence of the identification of inertia in IV Section 1, Proposition 1 and the subgroup compatibility in Proposition 2, using the local degree identity}
\uses{fw:inertia-relative-index,fw:local-inertia-size,fw:ramification-subgroups,fw:local-dvr-degree-identity,fw:fixed-field-integers}
\lean{FontaineWintenberger.inertia_fixedField_ramificationIdx, FontaineWintenberger.inertia_fixedField_top_ramificationIdx, FontaineWintenberger.inertia_fixedField_top_ramificationIdx_eq_finrank, FontaineWintenberger.inertia_fixedField_top_ramificationIdx_eq_field_finrank, FontaineWintenberger.inertia_fixedField_residue_surjective, FontaineWintenberger.inertia_fixedField_residueEquiv, FontaineWintenberger.inertia_fixedField_residueEquiv_apply}
\leanok
Let $E/K$ be a finite Galois extension under the standing complete
discrete-valuation and perfect-residue hypotheses. Put
$I=G_0\subseteq\operatorname{Gal}(E/K)$ and let $U=E^I$ be its actual
fixed intermediate field, with the restricted valuation. Then
\[
 e(U/K)=1,\qquad e(E/U)=[E:U],
\]
and the original residue inclusion $k_U\longrightarrow k_E$ is an
isomorphism. Thus $U/K$ is unramified and $E/U$ is totally ramified.
\end{lemma}
\begin{proof}
The inertia-relative index formula, applied to the subgroup $I$ itself,
gives $e(U/K)=[I:I]=1$. The zero lower group for the restricted
$I$-action on the same integer ring $\mathcal O_E$ is $I\cap I=I$,
by subgroup compatibility. The local inertia formula consequently gives
$e(E/U)=|I|$. The ordinary fixed-field theorem gives $[E:U]=|I|$.
Equivalently, Galois descent on the integer rings gives
$[\mathcal O_E:\mathcal O_U]=|I|$, where this degree denotes their
finite free module rank. The fundamental local degree identity over
$\mathcal O_U$ now gives $[k_E:k_U]=1$. The original residue map is
injective, since it is a unital homomorphism of fields, and degree one
makes it surjective; hence it is an isomorphism. Finally the residue
extension over $k_K$ is separable, because $k_K$ is perfect. Together
with $e(U/K)=1$ this is precisely unramifiedness of the finite extension
$U/K$, and the residue isomorphism is precisely total ramification of
$E/U$.
\end{proof}

\begin{lemma}[Inertia and the residue quotient]
\label{fw:inertia-quotient}
\dcref{Serre1979, Chapter IV, Section 1, Proposition 1 and its residue-quotient assertion}
\source{serre-1979-local-fields:page-0069}
\uses{fw:local-inertia-size,fw:inertia-fixed-field}
\lean{FontaineWintenberger.localResidueGaloisHom, FontaineWintenberger.localResidueGaloisHom_surjective, FontaineWintenberger.ker_localResidueGaloisHom, FontaineWintenberger.localResidueGaloisEquiv, FontaineWintenberger.spectralIntegerRing_residue_isGalois, FontaineWintenberger.spectralIntegerRing_card_lower_zero, FontaineWintenberger.inertia_fixedField_ramificationIdx, FontaineWintenberger.inertia_fixedField_top_ramificationIdx_eq_field_finrank, FontaineWintenberger.inertia_fixedField_residueEquiv}
\leanok
For finite Galois $E/K$, the residue extension is Galois, and reduction
gives an exact sequence
\[
 1\longrightarrow G_0\longrightarrow\Gal(E/K)
   \longrightarrow\Gal(k_E/k_K)\longrightarrow1.
\]
Thus $|G_0|=e(E/K)$. There is an unramified intermediate field $U/K$
with $k_U=k_E$, fixed by $G_0$, and $E/U$ is totally ramified.
\end{lemma}
\begin{proof}
The local residue-quotient theorem gives normality and separability of
$k_E/k_K$, surjectivity of reduction on the original automorphism
groups, and the exact zero lower group as its kernel. The same theorem
gives $|G_0|=e(E/K)$. Apply the preceding fixed-field lemma to the actual
field $U=E^{G_0}$: it gives $e(U/K)=1$, the original residue isomorphism
$k_U\simeq k_E$, and total ramification of $E/U$. By construction $G_0$
fixes $U$ pointwise, giving every assertion.
\end{proof}

\begin{lemma}[Computing inertia on a uniformizer]
\label{fw:uniformizer-ramification}
\dcref{Serre1979, Chapter IV, Section 1, reduction to the totally ramified case}
\source{serre-1979-local-fields:page-0070}
\uses{fw:inertia-fixed-field,fw:integral-generator,fw:inertia-quotient,fw:lower-ramification-properties,fw:ramification-subgroups}
\lean{FontaineWintenberger.ramificationOrder_eq_uniformizer_of_mem_inertia, FontaineWintenberger.ramificationOrder_eq_normalizedAdditiveValue_uniformizer, FontaineWintenberger.embeddingRamificationNumber_aut_eq_uniformizer_ratio, FontaineWintenberger.ramificationOrder_eq_zero_of_not_mem_inertia, FontaineWintenberger.embeddingRamificationNumber_aut_eq_neg_one_of_not_mem_inertia, FontaineWintenberger.finite_uniformizer_ratio_normalizedAdditiveValue_eq}
\leanok
For $s\in G_0$ and any uniformizer $\pi$ of $E$,
\[
 j_E(s)=v_E(s(\pi)-\pi),\qquad
 i_E(s)=v_E(s(\pi)/\pi-1).
\]
For $s\notin G_0$ one has $j_E(s)=0$ and $i_E(s)=-1$.
\end{lemma}
\begin{proof}
Let $U=E^{G_0}$. The extension $E/U$ is totally ramified and
$\mathcal O_E=\mathcal O_U[\pi]$. The generator formula over $U$ and
subgroup compatibility give the first equality. Subtracting
$v_E(\pi)=1$ gives the second. If $s\notin G_0$, some integral element
has nonzero residue difference, of valuation zero; all differences of
integral elements have nonnegative valuation. Their infimum is zero.
\end{proof}

\section{Quotient averaging and change of numbering}

\begin{lemma}[Tate's polynomial divisibility]
\label{fw:tate-polynomial-divisibility}
\dcref{Serre1979, Chapter IV, Section 1, Proposition 3, proof after J. Tate}
\source{serre-1979-local-fields:page-0070,serre-1979-local-fields:page-0071}
\uses{fw:generator-evaluation}
\lean{FontaineWintenberger.generator_sub_dvd_aut_sub, FontaineWintenberger.generator_sub_dvd_aeval_map, FontaineWintenberger.aeval_map_minpoly_dvd_generator_sub, FontaineWintenberger.associated_generator_sub_aeval_map_minpoly, FontaineWintenberger.addVal_generator_sub_eq_aeval_map_minpoly}
\leanok
Let $R\to A\hookrightarrow B$ be a compatible tower of commutative rings,
where $A$ is an integrally closed domain and $B$ is a domain. Suppose
$A=R[y]$, $B=R[x]$, and $x$ is integral over $A$. For an $R$-algebra
automorphism $\tau$ of $A$, put $f=\operatorname{minpol}_A(x)$ and let
$\tau(f)$ denote coefficientwise application of $\tau$. Then
\[
 a=\tau(y)-y,\qquad b=\tau(f)(x)
\]
are associated in $B$, where $a$ is mapped from $A$ into $B$.
In particular they have the same normalized additive valuation if $B$
is a DVR, including the value $+\infty$ at zero.
\end{lemma}
\begin{proof}
Every $c\in A=R[y]$ is a polynomial in $y$, so the difference
$\tau(c)-c$ is divisible by $\tau(y)-y$. Apply this to every coefficient
of $f$. The polynomial $\tau(f)-f$ is therefore divisible by the
constant polynomial $a$ after mapping its coefficients to $B$.
Evaluation at $x$, where $f(x)=0$, proves $a\mid b$.
Conversely write $y=g(x)$ with $g\in R[T]$, interpreting $y$ in $B$.
The polynomial $g(T)-y\in A[T]$ vanishes at $x$. The monic minimal
polynomial over the integrally closed domain $A$ divides every such
polynomial: this follows by passing to the fraction field, where
minimal-polynomial divisibility holds, and using monic division to
keep the quotient in $A[T]$. Hence $g(T)-y=f(T)q(T)$ for some
$q\in A[T]$. Applying $\tau$ to coefficients fixes those of $g$,
and evaluation at $x$ gives $y-\tau(y)=b\tau(q)(x)$.
Thus $b\mid a$ as well. Mutual divisibility in a domain is equivalent
to association. In a DVR a unit has additive valuation zero, so
associated elements have equal valuation, including zero.
\end{proof}

\begin{lemma}[The minimal polynomial of a Galois integral generator]
\label{fw:galois-generator-conjugates}
\dcref{Serre1979, Chapter IV, Section 1, Proposition 3}
\source{serre-1979-local-fields:page-0070,mathlib-invariant-basic,mathlib-is-galois-group}
\uses{fw:galois-ring-residue-theory}
\lean{FontaineWintenberger.map_minpoly_eq_prod_conjugates}
\leanok
Let a finite group $H$ act faithfully on a domain $B$, fixing exactly
an integrally closed subdomain $A$. If $B=A[x]$, then
\[
 \operatorname{minpol}_A(x)=\prod_{h\in H}(T-h(x))
\]
after mapping the coefficients from $A$ to $B$.
\end{lemma}
\begin{proof}
The product $P$ is monic of degree $|H|$, is fixed coefficientwise by
$H$, and vanishes at $x$ because of its identity factor. Its invariant
coefficients descend to $A$, producing a monic polynomial $p$ of the
same degree annihilating $x$. Thus $x$ is integral over $A$.
Since $A$ is integrally closed and $B=A[x]$ is a domain, the powers
$1,x,\ldots,x^{d-1}$ for $d=\deg\operatorname{minpol}_A(x)$ are an
$A$-basis of $B$: monic division spans, and linear independence follows
from minimal-polynomial divisibility over the fraction field.
Consequently $d=\operatorname{rank}_A B$.
The finite Galois ring action gives
$\operatorname{rank}_A B=|H|$, by passing to fraction fields and the
fixed-field degree theorem. The minimal polynomial divides $p$, as
$A$ is integrally closed. Both are monic of degree $|H|$, so they are
equal. Mapping to $B$ proves the formula.
\end{proof}

\begin{lemma}[The scaled Tate identity for integral rings]
\label{fw:scaled-tate-identity}
\dcref{Serre1979, Chapter IV, Section 1, Proposition 3}
\source{serre-1979-local-fields:page-0070,serre-1979-local-fields:page-0071}
\uses{fw:tate-polynomial-divisibility,fw:galois-generator-conjugates,fw:generator-evaluation,fw:ramification-index-scaling}
\lean{FontaineWintenberger.adjoin_singleton_eq_top_over_intermediate, FontaineWintenberger.map_minpoly_eq_prod_twisted, FontaineWintenberger.aeval_map_minpoly_eq_prod_displacements, FontaineWintenberger.addVal_aeval_map_minpoly_eq_sum, FontaineWintenberger.ramificationIdx_mul_order_eq_sum_coset}
\leanok
Let $R\to A\hookrightarrow B$ be a compatible tower with $A,B$ DVRs
and the last map local. Suppose $A=R[y]$, $B=R[x]$, and a finite
subgroup $H$ of a group $G$ acting by $R$-algebra automorphisms on
$B$ fixes exactly $A$. Let $\sigma\in G$ restrict to
$\tau\in\operatorname{Aut}_R(A)$. Then
\[
 e(B/A)\,j_A(\tau)=\sum_{h\in H}j_B(\sigma h),
\]
where each order is the intrinsic infimum of integral displacements.
The identity holds in $\mathbb N\cup\{+\infty\}$, including the
identity coset.
\end{lemma}
\begin{proof}
The generator $x$ also generates $B$ over $A$, since every polynomial
representation over $R$ is one over $A$. The finite Galois action
makes $x$ integral over $A$. The preceding factorization therefore
identifies $f=\operatorname{minpol}_A(x)$ with
$\prod_h(T-h(x))$. Compatibility of $\sigma$ and $\tau$ gives
\[
 \tau(f)(x)=\prod_{h\in H}(x-\sigma h(x)).
\]
Tate's divisibility lemma gives equality of its valuation with that of
$\tau(y)-y$. Multiplicativity of the additive valuation turns the
product into $\sum_h v_B(\sigma h(x)-x)$; changing signs leaves the
valuation unchanged. Scaling along $A\hookrightarrow B$ multiplies
$v_A(\tau(y)-y)$ by $e(B/A)$, and generator evaluation identifies
all these values with the stated intrinsic orders. This calculation
uses the extended nonnegative integers throughout, so a vanishing
factor gives $+\infty$ without an extra nonidentity assumption.
\end{proof}

\begin{lemma}[Compatible normed fields in an intermediate tower]
\label{fw:normed-intermediate-tower}
\dcref{Serre1979, Chapter II, Section 2, Proposition 3}
\source{serre-1979-local-fields:page-0035,mathlib-normed-valued}
\uses{fw:chosen-extension-integers}
\lean{NormedField.valuation_hasExtension, IntermediateField.normedAlgebraAmbient, IntermediateField.completeSpace_of_finiteDimensional, FontaineWintenberger.normedAlgebra_unitBall_isDiscreteValuationRing, FontaineWintenberger.normedAlgebra_unitBall_henselianLocalRing, FontaineWintenberger.normedAlgebra_unitBall_exists_generator}
\leanok
A normed algebra inclusion of nonarchimedean normed fields preserves
the norm and hence defines a compatible extension of their norm
valuations. For a complete base $K$ with DVR unit ball and a finite
separable normed algebra extension $F/K$, the actual unit ball of $F$
is finite over that of $K$, is a DVR, and is Henselian. If the base
residue field is perfect, it is monogenic over the base unit ball.
For a finite normed field extension $L/K$ and an intermediate field
$F$, the inherited norm makes $F$ a nontrivially normed field and
$L$ a normed algebra over $F$. The field $F$ is complete when $K$ is.
\end{lemma}
\begin{proof}
The normed-algebra bound at $1$ gives $|\iota(a)|\leq |a|$.
For $a\ne0$, apply the same bound to $a^{-1}$ and use multiplicativity
to obtain $|\iota(a)|^{-1}\leq |a|^{-1}$, hence the reverse inequality.
The assertion at zero is immediate. Multiplicativity now gives the
scalar norm identity, and the algebra map preserves the norm. Thus value comparisons agree under
restriction. The chosen-extension theorem applied to the actual norm
valuation of $F$ gives all stated integer-ring properties, including
monogenicity when the base residue is perfect.
For an intermediate field, the norm is the restriction of the ambient
norm. Its scalar action on $L$ is multiplication through this inclusion,
so the same norm identity proves the normed-algebra assertion. It is
nontrivial because it contains the nontrivially normed base. Finally
$F$ is a finite-dimensional normed vector space over the complete field
$K$, hence complete for this inherited norm; this is the finite-dimensional
completeness theorem with the original induced uniformity.
\end{proof}

\begin{lemma}[Actual integers in a finite Galois fixed field]
\label{fw:fixed-field-integers}
\dcref{Serre1979, Chapter II, Section 2; Chapter IV, Section 1, Proposition 3}
\source{serre-1979-local-fields:page-0035,serre-1979-local-fields:page-0070,mathlib-is-galois-group,mathlib-galois-group-degree}
\uses{fw:spectral-full-galois-action,fw:galois-ring-residue-theory,fw:chosen-extension-integers}
\lean{Valuation.HasExtension.comap_algebraMap, Valuation.HasExtension.valuationSubring_isScalarTower, IntermediateField.fixedField_isGaloisGroup, FontaineWintenberger.intermediate_spectral_hasExtension, FontaineWintenberger.intermediate_spectral_isDiscreteValuationRing, FontaineWintenberger.intermediate_spectral_henselianLocalRing, FontaineWintenberger.intermediate_spectral_exists_generator, FontaineWintenberger.fixedField_integer_moduleFinite, FontaineWintenberger.fixedField_integer_isGaloisGroup}
\leanok
Let $E/K$ be finite Galois over a complete nonarchimedean field with
DVR unit ball $R$. Put $S=\mathcal O_E$ for the actual spectral integer
ring. For a subgroup $H\subseteq\operatorname{Gal}(E/K)$ put $F=E^H$
and restrict the spectral valuation of $E$ to $F$, with actual
valuation ring $A$. Then $A$ is a DVR and is Henselian. Moreover $S$
is finite over $A$ and its $H$-invariants are exactly the image of $A$;
the induced action is faithful. If the residue field of $R$ is perfect,
then $A=R[y]$ and $S=R[x]$ for integral generators $y,x$.
\end{lemma}
\begin{proof}
Restrict the spectral valuation of $E$ along $F\hookrightarrow E$.
It extends the original norm valuation of $K$, since the composite
$K\to F\to E$ is the original inclusion. Apply the chosen-extension
integer theorem directly to $F/K$: its actual valuation ring $A$
is a finite DVR over $R$, and Henselian.
The inclusions $R\to A\to S$ are the restrictions of the field maps,
so form a scalar tower. A finite set spanning $S$ over $R$ also spans
it over $A$, and every element integral over $R$ is integral over $A$.
The fixed-field theorem identifies
$H$ as a faithful Galois group for $E/F$. Descent to the actual
integral rings, using integral closedness of the DVR $A$ and integrality
of $S/A$, therefore identifies the $H$-invariants of $S$ with
$A$, and proves faithfulness there. Finally apply the proved integral
generation theorem to the actual compatible extensions $F/K$ and
$E/K$, using perfection only for their finite residue extensions.
All valuation rings here are subrings of their stated fields and all
maps are their field inclusions.
\end{proof}

\begin{lemma}[Normal restriction on the actual fixed-field integers]
\label{fw:fixed-field-restriction}
\dcref{Serre1979, Chapter IV, Section 1, Proposition 3}
\source{serre-1979-local-fields:page-0070,mathlib-is-galois-group}
\uses{fw:fixed-field-integers,fw:spectral-full-galois-action,fw:generator-evaluation}
\lean{FontaineWintenberger.restrictNormal_mem_fixedField_decomposition, FontaineWintenberger.fixedFieldDecompositionRestriction, FontaineWintenberger.coe_fixedFieldDecompositionRestriction, FontaineWintenberger.coe_fixedFieldDecompositionRestriction_smul, FontaineWintenberger.fixedField_decompositionSubgroup_eq_top, FontaineWintenberger.fixedFieldDecompositionRestriction_surjective, FontaineWintenberger.fixedFieldDecompositionRestriction_ker, FontaineWintenberger.fixedFieldDecompositionQuotientEquiv, FontaineWintenberger.fixedFieldDecompositionQuotientEquiv_mk, FontaineWintenberger.fixedFieldRestrictionFiberEquiv, FontaineWintenberger.coe_fixedFieldRestrictionFiberEquiv, FontaineWintenberger.fixedFieldDecompositionRestriction_algebraMap, FontaineWintenberger.fixedFieldDecomposition_smulCommClass, FontaineWintenberger.fixedFieldDecomposition_order}
\leanok
In the preceding setting suppose $H$ is normal. Field restriction gives
a homomorphism $\rho:G=\operatorname{Gal}(E/K)\to\operatorname{Gal}(F/K)$,
and every $\rho(s)$ preserves the actual restricted valuation ring $A$.
Thus $\rho$ takes values in its decomposition group. For $a\in A$,
the inclusion $A\hookrightarrow S$ intertwines $\rho(s)(a)$ and $s(a)$.
The action on $A$ commutes with $R$, and its intrinsic ramification
order is the order of the induced $R$-algebra automorphism of $A$.
The map onto this decomposition group is surjective with kernel $H$.
In fact the decomposition group is all of $\operatorname{Gal}(F/K)$,
and restriction identifies it with $G/H$.
\end{lemma}
\begin{proof}
Normality of $H$ makes $F/K$ normal, so every $s\in G$ restricts to $F$.
For $a\in F$ the two images in $E$ satisfy
$\iota(\rho(s)(a))=s(\iota(a))$. The spectral norm is invariant under
$G$, so these images have the same norm. Consequently membership in
$A$ is preserved; applying the same argument to $s^{-1}$ gives equality
of valuation rings. Restriction respects products and the identity,
so gives the claimed homomorphism into the decomposition group.
Its intrinsic action is evaluation by this same field automorphism,
and the stated intertwining follows from the restriction identity.
It fixes the scalars in $R$ because the field restriction fixes $K$.
Finally both ramification orders are the infimum, over the same
$a\in A$, of the same normalized displacement valuations.
Every $K$-automorphism of $F$ extends to $E$, since $E/K$ is normal,
and its lift preserves $A$ by the norm argument above. Thus the
decomposition group is the full automorphism group, and restriction
onto it is surjective. The kernel consists of those elements fixing
$F$ pointwise, hence is $H$ by the finite Galois correspondence;
the first isomorphism theorem gives the quotient identification.
\end{proof}

\begin{lemma}[The coset averaging identity]
\label{fw:coset-ramification-average}
\dcref{Serre1979, Chapter IV, Section 1, Proposition 3, proof after J. Tate}
\source{serre-1979-local-fields:page-0070,serre-1979-local-fields:page-0071}
\uses{fw:integral-generator,fw:scaled-tate-identity,fw:fixed-field-integers,fw:fixed-field-restriction,fw:local-inertia-size,fw:integer-lower-groups}
\lean{FontaineWintenberger.fixedField_ramificationIdx_mul_order_eq_sum, FontaineWintenberger.fixedField_card_inertia, FontaineWintenberger.fixedField_card_inertia_mul_order_eq_sum, FontaineWintenberger.fixedField_order_eq_sum_div_ramificationIdx, FontaineWintenberger.fixedField_order_eq_sum_div_card_inertia, FontaineWintenberger.fixedField_order_eq_fiber_sum_div_ramificationIdx, FontaineWintenberger.fixedField_order_eq_fiber_sum_div_card_inertia}
\leanok
Let $H\triangleleft G=\Gal(E/K)$, $F=E^H$, and $Q=\Gal(F/K)=G/H$.
For $\bar s\in Q$,
\[
 j_F(\bar s)=\frac1{e(E/F)}\sum_{s\mapsto\bar s}j_E(s)
            =\frac1{|H_0|}\sum_{s\mapsto\bar s}j_E(s).
\]
For the identity coset the sum and both values are $+\infty$.
\end{lemma}
\begin{proof}
Only $\bar s\ne1$ needs proof. Choose integral generators $x$ of
$\mathcal O_E/\mathcal O_K$ and $y$ of
$\mathcal O_F/\mathcal O_K$, and a lift $s\in G$.
Write $f(T)=\prod_{t\in H}(T-t(x))\in\mathcal O_F[T]$ for the
minimal polynomial of $x$ over $F$. It is monic, so division of integral
polynomials by $f$ has integral quotient and remainder.
Set
\[
 a=s(y)-y,\qquad b=\prod_{t\in H}(st(x)-x).
\]
Every coefficient $c$ of $f$ lies in $\mathcal O_F=\mathcal O_K[y]$.
Consequently $s(c)-c$ is divisible by $a$, and evaluating
$s(f)-f$ at $x$ shows that $a$ divides $s(f)(x)=(-1)^{|H|}b$.

Conversely write $y=g(x)$ with $g\in\mathcal O_K[T]$.
Since $g(T)-y$ vanishes at $x$, division by $f$ gives
$g(T)-y=f(T)q(T)$ with $q\in\mathcal O_F[T]$.
Apply $s$ to the coefficients and evaluate at $x$ to get
$y-s(y)=s(f)(x)s(q)(x)$. Hence $b$ divides $a$.
The two nonzero elements have equal valuation. The generator formula now
gives
\[
 e(E/F)j_F(\bar s)=v_E(a)=v_E(b)
       =\sum_{t\in H}j_E(st).
\]
Finally $e(E/F)=|H_0|$ by the inertia lemma applied over $F$.
The kernel of restriction is $H$, so $t\mapsto st$ is a bijection from
$H$ onto the restriction fiber of $\bar s$. This reindexes the last sum
as the displayed fiber sum. Both denominators are positive finite
integers, so division in the extended nonnegative reals gives the
claimed average, also when the order is infinite.
\end{proof}

\begin{lemma}[Finite truncations of the intrinsic order]
\label{fw:truncated-order}
\dcref{Serre1979, Chapter IV, Section 3, Lemma 3}
\source{serre-1979-local-fields:page-0081}
\uses{fw:integral-order-function,fw:real-lower-groups}
\lean{FontaineWintenberger.truncatedRamificationOrder, FontaineWintenberger.truncatedRamificationOrder_of_order_top, FontaineWintenberger.truncatedRamificationOrder_of_order_nat, FontaineWintenberger.truncatedRamificationOrder_nonneg, FontaineWintenberger.truncatedRamificationOrder_le, FontaineWintenberger.continuous_truncatedRamificationOrder, FontaineWintenberger.monotone_truncatedRamificationOrder, FontaineWintenberger.truncatedRamificationOrder_one, FontaineWintenberger.truncatedRamificationOrder_zero, FontaineWintenberger.hasDerivAt_truncatedRamificationOrder, FontaineWintenberger.truncatedRamificationOrder_of_nonpos}
\leanok
For a group acting on a DVR, define
$T_g(t)=\min\{j(g),\max(t+1,0)\}$, viewed as a finite real number.
The cap is taken before conversion from extended to real numbers.
Each $T_g$ is continuous, increasing, nonnegative, and at most
$\max(t+1,0)$. If $j(g)=n<\infty$ it equals
$\min\{n,\max(t+1,0)\}$; if $j(g)=\infty$ it equals $\max(t+1,0)$.
In particular $T_1(t)=t+1$ for $t\geq-1$, and
$T_g(0)$ is the indicator of $g\in G_0$.
\end{lemma}
\begin{proof}
The finite cap makes the minimum finite even when the order is infinite.
For finite integer order the displayed formula follows directly from
the order-preserving inclusion of nonnegative integers into the reals;
for infinite order it follows from $\min(\infty,u)=u$.
Both formulas are continuous and increasing in $t$, since minimum with
a constant and positive part have these properties. The bounds are
immediate. The identity has infinite order. Finally the minimum of a
nonnegative integer (or infinity) with $1$ is $1$ exactly when that
order is at least $1$, which is precisely membership in $G_0$.
\end{proof}

\begin{lemma}[The integral equals the finite order sum]
\label{fw:herbrand-finite-sum}
\dcref{Serre1979, Chapter IV, Section 3, Lemma 3}
\source{serre-1979-local-fields:page-0080,serre-1979-local-fields:page-0081}
\uses{fw:real-lower-groups,fw:truncated-order,fw:integral-order-function}
\lean{FontaineWintenberger.antitone_card_lowerRamificationGroupReal, FontaineWintenberger.intervalIntegrable_card_lowerRamificationGroupReal, FontaineWintenberger.herbrandFunction, FontaineWintenberger.herbrandSum, FontaineWintenberger.sum_truncatedRamificationOrder_zero, FontaineWintenberger.herbrandSum_zero, FontaineWintenberger.continuous_herbrandSum, FontaineWintenberger.hasDerivAt_herbrandSum, FontaineWintenberger.integral_card_lower_eq_herbrandSum, FontaineWintenberger.herbrandSum_of_nonpos, FontaineWintenberger.herbrandFunction_eq_sum, FontaineWintenberger.continuousOn_herbrandFunction, FontaineWintenberger.hasDerivAt_herbrandFunction}
\leanok
For a finite group acting on a DVR, let $\varphi$ be the Herbrand function
defined by the actual lower-group cardinality integral, with
$\varphi(t)=t$ on $[-1,0]$. Then, for $t\geq-1$,
\[
 \varphi(t)=\frac{1}{|G_0|}\sum_{g\in G}T_g(t)-1.
\]
It is continuous on $[-1,\infty)$. For $t>-1$ outside the integer breaks
its derivative is $|G_t|/|G_0|$.
\end{lemma}
\begin{proof}
The lower groups decrease, so their finite cardinalities and the ratio
to the positive integer $|G_0|$ are decreasing functions. Hence that
ratio is integrable on each compact interval. The finite sum is continuous
by continuity of its capped summands. Outside the integer breaks a summand
has derivative $1$ if $t+1<j(g)$, and $0$ if $j(g)<t+1$.
The first condition is exactly $g\in G_t$ away from the breaks. An infinite
order contributes derivative $1$ as well. Thus differentiation of the
finite sum gives exactly the stated ratio. At zero its numerator is
$|G_0|$, by the preceding indicator identity, so the expression is zero.
The fundamental theorem of calculus, allowing the countable exceptional
set of integers, identifies the sum with the defining integral on $[0,t]$.
For $-1\leq t\leq0$, each summand is $t+1$ times its value at zero:
orders zero contribute zero and all orders at least one contribute $t+1$.
This proves equality also on $[-1,0]$. Continuity and the derivative
statement now transfer to the integral-defined function by this equality.
\end{proof}

\begin{lemma}[The Herbrand integral at an integer]
\label{fw:herbrand-integer-sum}
\uses{fw:herbrand-finite-sum}
\lean{FontaineWintenberger.herbrandFunction_nat_mul_card}
\leanok
For a finite group acting on a DVR and every integer $N\geq0$,
\[
 |G_0|\varphi_G(N)=\sum_{i=1}^{N}|G_i|.
\]
\end{lemma}
\begin{proof}
Both sides vanish at $N=0$. On the interval $(N,N+1]$ the real
lower group is $G_{N+1}$ by the ceiling convention. Integrating this constant cardinality on that interval gives
$|G_0|(\varphi_G(N+1)-\varphi_G(N))=|G_{N+1}|$.
Adding this equality to the formula at $N$ proves its successor,
and induction proves the assertion.
\end{proof}

\begin{lemma}[The Herbrand order isomorphism]
\label{fw:herbrand-order-isomorphism}
\dcref{Serre1979, Chapter IV, Section 3, Propositions 12--13}
\source{serre-1979-local-fields:page-0080,serre-1979-local-fields:page-0081}
\uses{fw:herbrand-finite-sum,fw:truncated-order}
\lean{FontaineWintenberger.strictMonoOn_herbrandSum, FontaineWintenberger.strictMonoOn_herbrandFunction, FontaineWintenberger.herbrandFunction_neg_one, FontaineWintenberger.herbrandFunction_mem_Ici, FontaineWintenberger.herbrandFunction_lower_bound, FontaineWintenberger.surjOn_herbrandFunction, FontaineWintenberger.herbrandOrderIso, FontaineWintenberger.herbrandOrderIso_apply, FontaineWintenberger.inverseHerbrandFunction, FontaineWintenberger.herbrandFunction_inverseHerbrandFunction, FontaineWintenberger.inverseHerbrandFunction_herbrandFunction, FontaineWintenberger.continuous_inverseHerbrandFunction, FontaineWintenberger.strictMonoOn_inverseHerbrandFunction, FontaineWintenberger.inverseHerbrandFunction_mem_Ici, FontaineWintenberger.herbrandFunction_of_nonpos, FontaineWintenberger.inverseHerbrandFunction_of_nonpos}
\leanok
The integral-defined Herbrand function is a strictly increasing continuous
bijection of $[-1,\infty)$ onto itself. Its inverse $\psi$ is strictly
increasing and continuous, with $\varphi(\psi(u))=u$ and
$\psi(\varphi(t))=t$ on that domain.
\end{lemma}
\begin{proof}
For $s<t$ with both indices at least $-1$, every capped summand is
increasing, and the identity summand increases strictly by $t-s$.
The finite sum therefore increases strictly. Its denominator $|G_0|$
is positive, so the same holds for $\varphi$ by the preceding equality.
At $-1$ its value is $-1$, hence its image lies in the same half-line.
Nonnegativity of all other summands gives the lower bound
$\varphi(t)\geq (t+1)/|G_0|-1$. For $u\geq-1$ choose
$T=|G_0|(u+1)-1\geq-1$; this bound gives $\varphi(T)\geq u$.
The intermediate value theorem on $[-1,T]$ supplies a preimage of $u$,
and strict monotonicity makes it unique. Thus $\varphi$ defines an order
isomorphism of the half-line. An order isomorphism of real intervals
is a homeomorphism for their order topologies. Its inverse is therefore
continuous and strictly increasing, and the two inverse identities follow.
\end{proof}

\begin{lemma}[Herbrand functions as finite sums]
\label{fw:herbrand-functions}
\dcref{Serre1979, Chapter IV, Section 3, Propositions 12--13 and Lemma 3}
\source{serre-1979-local-fields:page-0080,serre-1979-local-fields:page-0081}
\uses{fw:herbrand-finite-sum,fw:herbrand-order-isomorphism,fw:integral-order-function,fw:finite-valuation-ring,fw:spectral-full-galois-action}
\lean{FontaineWintenberger.truncatedRamificationOrder_affine, FontaineWintenberger.herbrandFunction_affine, FontaineWintenberger.herbrandFunction_affine_of_lower_eq_bot, FontaineWintenberger.exists_herbrandFunction_affine_tail, FontaineWintenberger.inverseHerbrandFunction_affine, FontaineWintenberger.exists_herbrand_affine_tails, FontaineWintenberger.herbrandFunction_finite_affine_partition, FontaineWintenberger.inverseHerbrandFunction_finite_affine_partition, FontaineWintenberger.exists_herbrandFunction_finite_derivative_breaks}
\leanok
For finite Galois $E/K$, $\varphi_{E/K}$ is a strictly increasing,
continuous, piecewise affine bijection of $[-1,\infty)$ onto itself,
and
\[
 \varphi_{E/K}(t)=\frac1{|G_0|}
      \sum_{s\in G}\min\{j_E(s),t+1\}-1.                 \tag{1}
\]
Its inverse $\psi_{E/K}$ is also continuous and piecewise affine.
Away from the finitely many breaks, for $t>0$,
$\varphi'_{E/K}(t)=|G_t|/|G_0|$.
\end{lemma}
\begin{proof}
Apply the general finite-action results to the faithful action of
$\operatorname{Gal}(E/K)$ on the actual DVR $\mathcal O_E$, supplied by
the finite-valuation-ring and full-Galois-action lemmas.
The finite-sum equality, continuity and derivative formula are proved in
Lemma~\ref{fw:herbrand-finite-sum}; the order isomorphism and continuous inverse
are proved in Lemma~\ref{fw:herbrand-order-isomorphism}.
For each nonnegative integer $n$ and $n\leq t\leq n+1$, a capped summand
is constant if $j(g)\leq n+1$, and has slope $1$ if $j(g)\geq n+2$,
including infinite order. Summing gives the explicit affine formula
\[
 \varphi(t)=\varphi(n)+(t-n)\frac{|G_{n+1}|}{|G_0|}.
\]
The slope is positive because each group contains the identity.
Eventual triviality supplies an integer $N\geq0$ above which every real
lower group is trivial. The derivative formula and the fundamental theorem
of calculus off the integer breaks then give
\[
 \varphi(t)=\varphi(N)+\frac{t-N}{|G_0|}\qquad(t\geq N).
\]
Together with $\varphi(t)=t$ on $[-1,0]$, the intervals with integer
endpoints from $0$ to $N$, and the final ray, form a finite affine
partition of the whole domain. Strict monotonicity takes these intervals
onto the corresponding intervals with endpoints $\varphi(n)$. Solving
an affine formula with positive slope gives
\[
 \psi(u)=n+\frac{u-\varphi(n)}{|G_{n+1}|/|G_0|}
 \quad\bigl(\varphi(n)\leq u\leq\varphi(n+1)\bigr),
\]
and $\psi(u)=N+(u-\varphi(N))|G_0|$ on the final ray.
Thus the inverse has a finite affine partition as well.
\end{proof}

\begin{lemma}[The first upper break of a finite fixing subgroup]
\label{fw:finite-first-upper-break}
\uses{fw:finite-upper-groups,fw:herbrand-functions,fw:real-lower-groups}
\lean{FontaineWintenberger.exists_first_upper_break_of_fixingSubgroup}
\leanok
Let a finite group $G$ act faithfully on a DVR. If a proper
subgroup $H$ satisfies $HG_0=G$, there is a finite real $b\geq0$
such that, for every $u\geq0$,
\[
 HG^u=G\quad\Longleftrightarrow\quad u\leq b.
\]
\end{lemma}
\begin{proof}
The lower groups are eventually trivial, while $H$ is proper.
Thus there is a least integer $n\geq0$ with $HG_n\ne G$.
The equality $HG_0=G$ implies $n\geq1$. Put $m=n-1$.
Monotonicity and minimality imply $HG_j=G$ exactly for $j<n$.
For $t\geq0$ the ceiling convention therefore gives
$HG_t=G$ exactly when $\lceil t\rceil\leq m$, equivalently
$t\leq m$. Put $b=\varphi_G(m)\geq0$. The inverse Herbrand
function is increasing and is the inverse of $\varphi_G$;
thus $\psi_G(u)\leq m$ exactly when $u\leq b$.
Since $G^u=G_{\psi_G(u)}$, this proves the assertion and includes
the endpoint $b$ without a continuity assumption on the filtration.
\end{proof}

\begin{lemma}[The largest order in a coset]
\label{fw:herbrand-coset-max}
\dcref{Serre1979, Chapter IV, Section 3, Lemma 4}
\source{serre-1979-local-fields:page-0081,serre-1979-local-fields:page-0082}
\uses{fw:coset-ramification-average,fw:herbrand-functions}
\lean{FontaineWintenberger.ramificationOrder_subgroup, FontaineWintenberger.ramificationOrder_mul_eq_min_of_coset_max, FontaineWintenberger.truncatedRamificationOrder_eq_coset_order, FontaineWintenberger.exists_fixedField_ramificationOrder_max, FontaineWintenberger.fixedField_order_sub_one_eq_herbrand, FontaineWintenberger.exists_fixedField_order_eq_herbrand}
\leanok
In the notation $H\triangleleft G$, $F=E^H$, $Q=G/H$, take
$\bar s\ne1$ and put $m=\max_{s\mapsto\bar s}j_E(s)$. Then
\[
 j_F(\bar s)-1=\varphi_{E/F}(m-1).
\]
\end{lemma}
\begin{proof}
The restriction fiber is finite and nonempty. Every element in it is
nonidentity, since its restriction is nonidentity; faithfulness therefore
makes every order finite. Choose a lift $s$ with $j_E(s)=m$.
For $t\in H$ with $j_E(t)<m$,
the strict product rule gives $j_E(st)=j_E(t)$.
For $j_E(t)\geq m$ the product inequality and maximality of $m$ give
$j_E(st)=m$. Thus $j_E(st)=\min\{j_E(t),m\}$ in both cases.
The averaging identity, subgroup compatibility, and (1) for $E/F$ imply
\[
 j_F(\bar s)-1=
 \frac1{|H_0|}\sum_{t\in H}\min\{j_E(t),m\}-1
 =\varphi_{E/F}(m-1).
\]
This also covers $m=0$, since $\varphi_{E/F}(-1)=-1$.
\end{proof}

\begin{theorem}[Herbrand's lower quotient formula]
\label{fw:herbrand-lower-quotients}
\dcref{Serre1979, Chapter IV, Section 3, Lemma 5}
\source{serre-1979-local-fields:page-0082}
\uses{fw:herbrand-coset-max}
\lean{FontaineWintenberger.mem_lowerRamificationGroupReal_iff_le_toNat, FontaineWintenberger.fixedField_map_lowerRamificationGroupReal, FontaineWintenberger.fixedField_quotient_lowerRamificationGroupReal}
\leanok
With $H\triangleleft G$ and $F=E^H$, for every $t\geq-1$,
\[
 G_tH/H=(G/H)_{\varphi_{E/F}(t)}.
\]
\end{theorem}
\begin{proof}
The identity belongs to both sides. For a nonidentity coset $\bar s$,
with maximum $m$ as above,
\[
 \bar s\in G_tH/H
 \ \Longleftrightarrow\ m-1\geq t
 \ \Longleftrightarrow\ \varphi_{E/F}(m-1)\geq\varphi_{E/F}(t)
 \ \Longleftrightarrow\ j_F(\bar s)-1\geq\varphi_{E/F}(t).
\]
The last condition is exactly membership in the right side. The
equivalence also holds at $t=-1$ and at the integral endpoints because
the groups use ceilings and weak inequalities.
\end{proof}

\begin{theorem}[Transitivity of Herbrand functions]
\label{fw:herbrand-transitivity}
\dcref{Serre1979, Chapter IV, Section 3, Proposition 15}
\source{serre-1979-local-fields:page-0081,serre-1979-local-fields:page-0082}
\uses{fw:herbrand-lower-quotients,fw:ramification-subgroups,fw:herbrand-functions}
\lean{FontaineWintenberger.fixedField_card_lowerRamificationGroupReal, FontaineWintenberger.fixedField_card_lowerRamificationGroup_zero, FontaineWintenberger.fixedField_herbrand_slope_mul, FontaineWintenberger.fixedField_exists_herbrand_comp_derivative_breaks, FontaineWintenberger.fixedField_herbrandFunction_transitivity, FontaineWintenberger.fixedField_inverseHerbrandFunction_transitivity}
\leanok
If $E/K$ and $F/K$ are finite Galois and $K\subset F\subset E$, then
\[
 \varphi_{E/K}=\varphi_{F/K}\circ\varphi_{E/F},\qquad
 \psi_{E/K}=\psi_{E/F}\circ\psi_{F/K}.
\]
\end{theorem}
\begin{proof}
Put $G=\Gal(E/K)$, $H=\Gal(E/F)$, $Q=G/H$, and
$v=\varphi_{E/F}(t)$. For $t\geq0$, lower quotient compatibility and
subgroup compatibility give an exact sequence of finite groups
\[
 1\longrightarrow H_t\longrightarrow G_t\longrightarrow Q_v
 \longrightarrow1.
\]
At zero this gives $|G_0|=|H_0||Q_0|$; for any $t$ it gives
$|G_t|=|H_t||Q_v|$.
Outside the finitely many breaks of $\varphi_{E/F}$ and the inverse
images of the breaks of $\varphi_{F/K}$, the chain rule gives
\[
 (\varphi_{F/K}\circ\varphi_{E/F})'(t)
    =\frac{|Q_v|}{|Q_0|}\frac{|H_t|}{|H_0|}
    =\frac{|G_t|}{|G_0|}=\varphi'_{E/K}(t).
\]
The exceptional set is finite: inverse images of breaks of
$\varphi_{F/K}$ are obtained by applying $\psi_{E/F}$ to those breaks.
The actual function $|G_t|/|G_0|$ is integrable on compact intervals by
monotonicity. The fundamental theorem of calculus off this finite set
therefore identifies the composite, which vanishes at zero, with its
integral from zero. That integral is $\varphi_{E/K}$ by definition.
On $[-1,0]$ all three functions are the identity.
Taking inverses proves the second identity.
\end{proof}

\begin{theorem}[Herbrand compatibility with quotients]
\label{fw:herbrand-quotients}
\dcref{Serre1979, Chapter IV, Section 3, Proposition 14}
\source{serre-1979-local-fields:page-0081,serre-1979-local-fields:page-0082}
\uses{fw:herbrand-transitivity,fw:herbrand-lower-quotients,fw:finite-upper-groups}
\lean{FontaineWintenberger.fixedField_map_upperRamificationGroup, FontaineWintenberger.fixedField_quotient_upperRamificationGroup}
\leanok
If $E/K$ is finite Galois and $F/K$ is a Galois intermediate extension,
restriction maps $\Gal(E/K)^u$ onto $\Gal(F/K)^u$ for every $u\geq-1$.
\end{theorem}
\begin{proof}
Put $H=\Gal(E/F)$, so $E^H=F$ by finite Galois correspondence.
For $t=\psi_{E/K}(u)$, transitivity gives
$\varphi_{E/F}(t)=\psi_{F/K}(u)$. Lower quotient compatibility therefore
identifies the image of $G_t=G^u$ with
$\Gal(F/K)_{\psi_{F/K}(u)}=\Gal(F/K)^u$.
\end{proof}

\begin{corollary}[Quotients by a lower ramification group]
\label{fw:ramification-fixed-quotient}
\dcref{Serre1979, Chapter IV, Section 1, Corollary to Proposition 3; equivalently, its lower-Herbrand quotient form}
\source{serre-1979-local-fields:page-0071}
\uses{fw:herbrand-lower-quotients,fw:herbrand-functions,fw:lower-ramification-properties,fw:ramification-subgroups}
\lean{FontaineWintenberger.lowerFixedField_quotient_lowerRamificationGroupReal,FontaineWintenberger.lowerFixedField_lowerRamificationGroupReal_eq_bot}
\leanok
Let $r\geq0$ be an integer, $H=G_r$, and $F=E^H$.
Under the canonical restriction identification
$G/H\simeq\operatorname{Gal}(F/K)$, for every real $t$ with
$-1\leq t\leq r$,
\[
 (G/H)_t=G_t/H,
\]
where the right side denotes the canonical image of $G_t$.
For every real $t\geq r$, $(G/H)_t=\{1\}$.
In particular both formulas apply at $t=r$ and also when $r=0$.
\end{corollary}
\begin{proof}
The lower group $H=G_r$ is normal and is contained in $G_0$.
Subgroup compatibility identifies its relative lower groups with
$H_s=H\cap G_s$. For $-1\leq s\leq r$, monotonicity gives
$H\subseteq G_s$, and hence $H_s=H$. In particular $H_0=H$.
Its relative Herbrand function satisfies
\[
 \varphi_{E/F}(t)=t\qquad(-1\leq t\leq r).
\]
For $-1\leq t\leq0$ this is its definition. For $0\leq t\leq r$,
the integrand $|H_s|/|H_0|$ is identically one on $[0,t]$, so the
defining integral is $t$. The proved lower quotient formula therefore
gives
\[
 \rho(G_t)=(G/H)_{\varphi_{E/F}(t)}=(G/H)_t
 \qquad(-1\leq t\leq r),
\]
where $\rho:G\to\operatorname{Gal}(F/K)$ is the original restriction.
Since $H\subseteq G_t$ in this range, its image is exactly $G_t/H$.
At $t=r$, this is the image of $H$, which is the identity subgroup
because $\ker\rho=H$. The quotient lower filtration is decreasing,
so it stays the identity subgroup for every $t\geq r$.
\end{proof}

\section{Absolute groups and finite-level consumers}

\begin{lemma}[Field isomorphisms preserve intrinsic ramification]
\label{fw:ramification-field-equiv}
\dcref{Serre1979, Chapter IV, Sections 1 and 3, definitions}
\source{serre-1979-local-fields:page-0068,serre-1979-local-fields:page-0080}
\uses{fw:spectral-integers,fw:finite-valuation-ring,fw:ramification-dvr-equiv,fw:spectral-full-galois-action}
\lean{FontaineWintenberger.spectralIntegerRingEquiv, FontaineWintenberger.coe_spectralIntegerRingEquiv, FontaineWintenberger.spectralIntegerRingEquiv_smul, FontaineWintenberger.spectralIntegerRingEquiv_addVal, FontaineWintenberger.ramificationOrder_autCongr, FontaineWintenberger.herbrandFunction_autCongr, FontaineWintenberger.upperRamificationGroup_map_autCongr}
\leanok
For a $K$-isomorphism $f:E\simeq E'$ of finite Galois extensions,
$f$ restricts to an isomorphism of the actual integer rings. Conjugation
$s\mapsto f s f^{-1}$ preserves intrinsic orders, Herbrand functions
and upper ramification groups.
\end{lemma}
\begin{proof}
The actual integer rings are the integral closures of $\mathcal O_K$.
An algebra isomorphism maps integral elements onto integral elements;
thus the induced integral-closure isomorphism restricts the given field
map on every element. It is equivariant for conjugation, since
$f(s(a))=(f s f^{-1})(f(a))$. The preceding DVR-isomorphism lemma
then identifies normalized valuations, intrinsic orders, real lower
groups, their defining Herbrand integrals and inverse upper numbering.
\end{proof}

\begin{lemma}[Ordinary restriction on finite Galois upper groups]
\label{fw:ordinary-upper-restriction}
\dcref{Serre1979, Chapter IV, Section 3, Proposition 14}
\source{serre-1979-local-fields:page-0081,serre-1979-local-fields:page-0082}
\uses{fw:herbrand-quotients,fw:ramification-action-equiv,fw:spectral-valuation-ring-unique}
\lean{FontaineWintenberger.restrictNormal_mem_intermediate_decomposition, FontaineWintenberger.intermediateDecompositionRestriction, FontaineWintenberger.intermediate_map_upperRamificationGroup, FontaineWintenberger.restrictNormalHom_map_upperRamificationGroup}
\leanok
Let $E/K$ be finite Galois and $F\subset E$ a Galois intermediate field.
With the actual extended valuation rings on both fields, ordinary field
restriction maps $\Gal(E/K)^u$ onto $\Gal(F/K)^u$ for every $u\geq-1$.
\end{lemma}
\begin{proof}
Put $H=\Gal(E/F)$; finite Galois correspondence identifies $E^H$ with
$F$. Restriction into the decomposition group of the pullback of
$\mathcal O_E$ to $F$ maps upper groups onto upper groups by Herbrand's
quotient theorem. Uniqueness of the extension valuation ring identifies
that pullback with the actual ring $\mathcal O_F$. Its decomposition
group is the full group $\Gal(F/K)$, and the resulting equivalence
identifies the actions on every integral element. Equivariant relabelling
therefore identifies the upper groups. The composite restriction and
inclusion is the ordinary field restriction on every field element,
which proves the assertion for that map.
\end{proof}

\begin{lemma}[Restriction for finite Galois field towers]
\label{fw:finite-tower-upper-restriction}
\dcref{Serre1979, Chapter IV, Section 3, Proposition 14}
\source{serre-1979-local-fields:page-0081,serre-1979-local-fields:page-0082}
\uses{fw:ordinary-upper-restriction,fw:ramification-field-equiv}
\lean{FontaineWintenberger.restrictNormalHom_map_upperRamificationGroup_tower, FontaineWintenberger.finGaloisGroupMap_map_upperRamificationGroup}
\leanok
In any compatible field tower $K\to F\to E$ with $F/K$ and $E/K$
finite Galois, ordinary field restriction maps $\Gal(E/K)^u$ onto
$\Gal(F/K)^u$, using each field's own integer ring.
In particular the ordinary transition maps between finite Galois
subextensions of a common field preserve these actual upper groups.
\end{lemma}
\begin{proof}
Let $F'$ be the image of $F$ in $E$, and let $f:F\simeq F'$ be the
canonical range isomorphism. Its induced conjugation equivalence maps
the actual upper groups of $F$ onto those of $F'$. Restriction into
$F'$ maps the actual upper groups of $E$ onto those of $F'$ by the
intermediate-field assertion. Conjugation by $f$ composed with restriction
to $F$ equals restriction to $F'$: after embedding into $E$, both maps
send each $a\in F$ to $s(a)$. Injectivity of the conjugation equivalence
then identifies the two subgroup images and proves the assertion.
For finite Galois subextensions of a common field, the transition map is
ordinary restriction for their canonical inclusion tower, so the same
identity applies.
\end{proof}

\begin{lemma}[Profinite Galois groups and compact intersections]
\label{fw:profinite-galois-foundations}
\dcref{MathlibProfiniteGalois, Profinite, Normal, and Compactness modules}
\source{mathlib-profinite-galois,mathlib-normal-basic,mathlib-compact-finite-intersection}
\lean{InfiniteGalois.instCompactSpaceAlgEquivOfIsGalois, InfiniteGalois.continuousMulEquivToLimit, InfiniteGalois.restrictNormalHom_continuous, AlgEquiv.restrictNormalHom_surjective, IsCompact.inter_iInter_nonempty}
\mathlibok
For an algebraic Galois extension $\Omega/K$, its actual Galois group,
with its Krull topology, is compact and is canonically isomorphic as a
topological group to the inverse limit of its finite Galois groups.
Restriction to a normal intermediate field is continuous and surjective.
If $C$ is compact and $(T_i)$ is a family of closed sets such that
$C\cap\bigcap_{i\in S}T_i$ is nonempty for every finite $S$, then
$C\cap\bigcap_iT_i$ is nonempty.
\end{lemma}
\begin{proof}
The restriction maps are compatible. Conversely a compatible tuple
$(q_E)_E$ defines an automorphism by $q(x)=q_E(x)$ in any finite Galois
subextension containing $x$. Such a subextension exists by taking the
normal closure of $K(x)/K$. Two choices lie in their finite Galois
compositum, so compatibility proves independence. Put any finitely many
elements in one finite Galois subextension to check addition,
multiplication and base-field linearity. The tuple of inverses defines
the inverse automorphism. These constructions are mutually inverse,
and the finite fixing-subgroup bases of the Krull topologies show
continuity in both directions. The inverse limit is a closed subset
of the product of the finite discrete Galois groups: each compatibility
condition is a closed equality. The product, and hence the limit, is
compact. Restriction is continuous on the same finite fixing-subgroup
bases and is surjective by the normal-extension automorphism extension
theorem. For the last assertion, an empty intersection would make the
open complements of the $T_i$ cover $C$. A finite subcover would
contradict the assumed finite intersection property.
\end{proof}

\begin{lemma}[Infinite upper ramification groups]
\label{fw:infinite-upper-groups}
\dcref{Serre1979, Chapter IV, Section 3, Remark 1 after Proposition 14}
\source{serre-1979-local-fields:page-0082}
\uses{fw:finite-upper-groups,fw:finite-valuation-ring,fw:spectral-full-galois-action,fw:profinite-galois-foundations}
\lean{FontaineWintenberger.upperRamificationGroupInfinite, FontaineWintenberger.mem_upperRamificationGroupInfinite, FontaineWintenberger.upperRamificationGroupInfinite_normal, FontaineWintenberger.upperRamificationGroupInfinite_antitoneOn, FontaineWintenberger.upperRamificationGroupInfinite_neg_one, FontaineWintenberger.isClosed_upperRamificationGroupInfinite, FontaineWintenberger.map_upperRamificationGroupInfinite_le, FontaineWintenberger.isCompact_upperRamificationGroupInfinite}
\leanok
For an algebraic Galois extension $\Omega/K$, put
\[
 \Gamma_{\Omega/K}^u=
 \bigcap_{E\subset\Omega,\ E/K\ \mathrm{finite\ Galois}}
   \operatorname{res}_{E}^{-1}(\Gal(E/K)^u).
\]
These are closed normal subgroups of the actual Krull-topological group
$\Gal(\Omega/K)$, decreasing for $u\geq-1$, and
$\Gamma_{\Omega/K}^{-1}=\Gal(\Omega/K)$.
They are compact.
\end{lemma}
\begin{proof}
Every finite upper group is a normal subgroup, and preimages and
intersections preserve normality. Restriction is continuous for the
Krull topologies, and a subgroup of a finite discrete Galois group is
closed. The displayed intersection is therefore closed. Decreasingness
and the value at $-1$ follow on every finite level.
Closedness in the compact Galois group also proves compactness.
\end{proof}

\begin{lemma}[Surjective finite upper projections]
\label{fw:finite-upper-projections}
\dcref{Serre1979, Chapter IV, Section 3, Remark 1 after Proposition 14}
\source{serre-1979-local-fields:page-0082}
\uses{fw:infinite-upper-groups,fw:finite-tower-upper-restriction,fw:profinite-galois-foundations}
\lean{FontaineWintenberger.exists_lift_finite_upper_conditions, FontaineWintenberger.exists_lift_upperRamificationGroupInfinite, FontaineWintenberger.map_upperRamificationGroupInfinite}
\leanok
For every finite Galois $E\subset\Omega$ and $u\geq-1$, ordinary
restriction maps $\Gamma_{\Omega/K}^u$ onto $\Gal(E/K)^u$.
\end{lemma}
\begin{proof}
Fix $q\in\Gal(E/K)^u$. Its restriction fiber $C$ in
$\Gal(\Omega/K)$ is closed and hence compact. For every finite Galois
$D\subset\Omega$, let $T_D$ be the closed condition that restriction
lies in $\Gal(D/K)^u$. Given finitely many $D$, form their compositum
$M$ together with $E$. Upper quotient compatibility lifts $q$ to some
$r\in\Gal(M/K)^u$. Extend $r$ to $s\in\Gal(\Omega/K)$ by normality.
Compatibility of ordinary restriction shows $s|_E=q$ and
$s|_D\in\Gal(D/K)^u$ for every chosen $D$. Thus $C$ and the $T_D$
have the finite intersection property. Compactness gives a point in
$C\cap\bigcap_D T_D$, which is the required lift in
$\Gamma_{\Omega/K}^u$. The reverse inclusion of subgroup images is
the defining finite-level membership condition.
\end{proof}

\begin{lemma}[Upper images in Galois subextensions]
\label{fw:infinite-upper-restriction}
\dcref{Serre1979, Chapter IV, Section 3, Remark 1 after Proposition 14}
\source{serre-1979-local-fields:page-0082}
\uses{fw:infinite-upper-groups,fw:finite-upper-projections,fw:ramification-field-equiv,fw:profinite-galois-foundations}
\lean{FontaineWintenberger.finiteGaloisLift, FontaineWintenberger.finiteGaloisLiftEquiv, FontaineWintenberger.restrictNormal_finiteGaloisLift, FontaineWintenberger.restrictNormal_mem_upperRamificationGroupInfinite, FontaineWintenberger.exists_lift_finite_infinite_upper_restrictions, FontaineWintenberger.exists_lift_infinite_upper_restriction, FontaineWintenberger.restrictNormalHom_map_upperRamificationGroupInfinite}
\leanok
Let $L$ be a Galois intermediate field of an algebraic Galois extension
$\Omega/K$. For every $u\geq-1$, ordinary restriction satisfies
\[
 \operatorname{res}_{L}(\Gamma_{\Omega/K}^u)=\Gamma_{L/K}^u.
\]
\end{lemma}
\begin{proof}
A finite Galois subfield $E$ of $L$ has an actual image $\widetilde E$
in $\Omega$, canonically isomorphic to $E$ over $K$. On automorphism
groups, restriction to $\widetilde E$ is conjugation of the composite
restriction through $L$ to $E$. The isomorphism theorem for upper groups
therefore gives the inclusion from left to right.

Fix $q\in\Gamma_{L/K}^u$. Inside the compact group
$\Gamma_{\Omega/K}^u$, impose, for each finite Galois $E\subset L$, the
closed condition $(s|_L)|_E=q|_E$. These conditions have the finite
intersection property. Indeed let $E_0$ be the compositum inside $L$
of a finite list of such $E$. Its restriction $q|_{E_0}$ belongs to
$\Gal(E_0/K)^u$. Conjugate along $E_0\simeq\widetilde E_0$ and apply
surjective finite upper projection to lift it to
$s\in\Gamma_{\Omega/K}^u$. The conjugation identity and compatibility
of restrictions show that this $s$ satisfies every condition in the list.
Compactness gives an $s$ satisfying all finite conditions. The canonical
Galois inverse-limit isomorphism for $L/K$ is injective, so equality on
all finite Galois subextensions implies $s|_L=q$. This proves the reverse
inclusion.
\end{proof}

\begin{lemma}[The actual inverse limit of finite upper groups]
\label{fw:upper-inverse-limit}
\dcref{Serre1979, Chapter IV, Section 3, Remark 1 after Proposition 14}
\source{serre-1979-local-fields:page-0082}
\uses{fw:infinite-upper-groups,fw:finite-tower-upper-restriction,fw:profinite-galois-foundations}
\lean{FontaineWintenberger.finiteUpperRamificationGroupMap, FontaineWintenberger.coe_finiteUpperRamificationGroupMap, FontaineWintenberger.finiteUpperRamificationGroupMap_id, FontaineWintenberger.finiteUpperRamificationGroupMap_comp, FontaineWintenberger.finiteUpperRamificationGroupFunctor, FontaineWintenberger.upperRamificationToLimit, FontaineWintenberger.upperRamificationLimitForget, FontaineWintenberger.upperRamificationLimitForget_toLimit, FontaineWintenberger.continuous_upperRamificationToLimit, FontaineWintenberger.upperRamificationToLimit_bijective, FontaineWintenberger.upperRamificationEquivToLimit, FontaineWintenberger.upperRamificationEquivToLimit_apply_val}
\leanok
For an algebraic Galois extension $\Omega/K$ and $u\geq-1$, the actual
finite upper groups, with ordinary restriction maps, form an inverse
system of finite topological groups. The map
\[
 \Gamma_{\Omega/K}^u\longrightarrow
 \varprojlim_{E\subset\Omega,\ E/K\ \mathrm{finite\ Galois}}\Gal(E/K)^u,
 \qquad s\longmapsto(s|_E)_E
\]
is an isomorphism of topological groups, with each coordinate equal to
ordinary restriction.
\end{lemma}
\begin{proof}
The finite upper quotient theorem shows that each ordinary restriction
maps the source upper group into the target upper group. Identity and
composition follow from the same properties of ordinary restriction;
all finite groups carry their inherited discrete Krull topology.
An upper automorphism gives a compatible tuple by restriction. Forgetting
upper membership gives a tuple in the canonical inverse limit of the
full finite Galois groups. Its actual Galois automorphism, supplied by
the preceding canonical isomorphism, belongs to
$\Gamma_{\Omega/K}^u$ because every coordinate lies in its prescribed
upper subgroup. This is inverse to restriction, proving bijectivity
and the group law. The forward map is continuous since all its finite
coordinates are continuous restrictions. Its domain is compact and its
target is Hausdorff, so this continuous bijection is a homeomorphism.
\end{proof}

\begin{proposition}[Surjective absolute upper projections]
\label{fw:absolute-upper-projections}
\dcref{Serre1979, Chapter IV, Section 3, Remark 1 after Proposition 14}
\source{serre-1979-local-fields:page-0082}
\uses{fw:herbrand-quotients,fw:ordinary-upper-restriction,fw:finite-tower-upper-restriction,fw:infinite-upper-groups,fw:finite-upper-projections,fw:infinite-upper-restriction,fw:upper-inverse-limit}
\lean{FontaineWintenberger.absoluteUpperRamificationGroup, FontaineWintenberger.absolute_map_upperRamificationGroup, FontaineWintenberger.absolute_restrict_map_upperRamificationGroupInfinite, FontaineWintenberger.absoluteUpperImageEquivToLimit, FontaineWintenberger.absoluteUpperImageEquivToLimit_apply_val, FontaineWintenberger.upperRamificationGroupInfinite_normal, FontaineWintenberger.upperRamificationGroupInfinite_antitoneOn, FontaineWintenberger.isClosed_upperRamificationGroupInfinite, FontaineWintenberger.isCompact_upperRamificationGroupInfinite}
\leanok
For $u\geq-1$ and every finite Galois $E/K$, restriction maps $G_K^u$ onto
$\Gal(E/K)^u$. More generally, for any Galois algebraic extension $L/K$,
the image $\Gamma^u$ of $G_K^u$ in $\Gamma=\Gal(L/K)$ is
\[
 \Gamma^u\simeq\varprojlim_{E\subset L,\ E/K\ \mathrm{finite\ Galois}}
                   \Gal(E/K)^u.
\]
The isomorphism sends an automorphism to its ordinary finite restrictions.
Here $G_K^u=\Gamma_{K^{\mathrm{sep}}/K}^u$ inside the actual absolute
Galois group. Each $G_K^u$ and $\Gamma^u$ is closed and normal, and these groups
decrease with $u$.
\end{proposition}
\begin{proof}
Fix $q\in\Gal(E/K)^u$. In the compact profinite group $G_K$, impose
the closed conditions $g|_E=q$ and
$g|_D\in\Gal(D/K)^u$ for every finite Galois $D/K$.
Any finite set of these conditions is satisfiable: take the compositum
$M$ of $E$ and the finitely many $D$. The finite quotient theorem lifts
$q$ to $\Gal(M/K)^u$, whose restriction satisfies all the other
conditions, and a finite Galois automorphism extends to $G_K$.
Compactness gives an element satisfying every condition. It lies in
$G_K^u$ and restricts to $q$.

For the second assertion let $(q_E)_E$ be an element of the displayed
inverse limit. Impose $g|_E=q_E$ for every finite Galois $E/K$ in $L$,
together with all the upper membership conditions for finite Galois
$D/K$ in $K^{\mathrm{sep}}$. For finitely many such conditions, let
$E_0$ be the compositum of the finitely many $E$, inside $L$, and let
$M$ also include the finitely many $D$. Lift $q_{E_0}$ to
$\Gal(M/K)^u$ by the finite quotient theorem and then extend to $G_K$.
Compactness again supplies a common solution. This proves surjectivity
onto the inverse limit. The canonical inverse-limit identification for
$L/K$ proves injectivity on the image, and its restriction to the closed
upper group is a topological group isomorphism by the preceding lemma.
Normality and decreasingness hold on every finite level. Closedness of
$G_K^u$ follows from its definition by closed conditions, and closedness
of its image follows from compactness and the Hausdorff property of
$\Gamma$.
\end{proof}

\begin{lemma}[The inverse Herbrand integral]
\label{fw:upper-inverse-integral}
\dcref{Serre1979, Chapter IV, Section 3, formula following Proposition 13}
\source{serre-1979-local-fields:page-0081}
\uses{fw:herbrand-functions,fw:finite-upper-groups}
\lean{FontaineWintenberger.upperRamificationGroup_relIndex_eq_div, FontaineWintenberger.intervalIntegrable_inverseHerbrand_slope, FontaineWintenberger.exists_inverseHerbrandFunction_finite_derivative_breaks, FontaineWintenberger.inverseHerbrandFunction_eq_integral_div, FontaineWintenberger.inverseHerbrandFunction_eq_integral_relIndex}
\leanok
For finite Galois $E/K$ and $u\geq0$,
\[
 \psi_{E/K}(u)=\int_0^u[G^0:G^v]\,dv.
\]
\end{lemma}
\begin{proof}
On an open affine piece, put $t=\psi_{E/K}(v)$.
The inverse slope is
$\psi'_{E/K}(v)=|G_0|/|G_t|=[G^0:G^v]$.
The inverse derivative formula follows by differentiating the proved local
identity $\varphi_{E/K}(\psi_{E/K}(v))=v$; its derivative is invertible
because every lower group contains the identity. The exceptional values
are among the images under $\varphi_{E/K}$ of the finitely many breaks.
For $v\geq0$, decreasingness gives $G^v\subset G^0$, and Lagrange's
formula gives the displayed index equality. The cardinality ratio is
increasing on $[0,u]$, hence integrable there. Apply the fundamental
theorem of calculus off the finite exceptional set, using continuity and
$\psi_{E/K}(0)=0$, to obtain the formula including endpoints.
\end{proof}


\begin{lemma}[The Herbrand integral for an arbitrary subgroup]
\label{fw:herbrand-subgroup-integral}
\dcref{Wintenberger1983, Proposition 1.2.3; finite subgroup change of parameter}
\source{wintenberger:page-0005,serre-1979-local-fields:page-0081}
\uses{fw:upper-inverse-integral,fw:ramification-subgroups,fw:herbrand-functions}
\lean{FontaineWintenberger.card_lowerRamificationGroupReal_subgroup, FontaineWintenberger.card_lowerRamificationGroup_zero_subgroup, FontaineWintenberger.subgroup_upperRelIndex_mul_card, FontaineWintenberger.subgroup_herbrand_slope_mul, FontaineWintenberger.subgroup_upperRelIndex_pos, FontaineWintenberger.subgroup_upperRelIndex_monotoneOn, FontaineWintenberger.intervalIntegrable_subgroup_upperRelIndex, FontaineWintenberger.exists_subgroup_herbrand_comp_derivative_breaks, FontaineWintenberger.subgroup_herbrandFunction_comp_inverse_eq_integral, FontaineWintenberger.subgroup_herbrandFunction_comp_inverse_of_nonpos}
\leanok
Let a finite group $G$ act faithfully on a discrete valuation ring, let
$H\subset G$ be any subgroup, and use the actual lower and upper groups
of these actions. Write $I=G_0$, $A=H\cap I$, and, for $u\geq0$, put
\[
 b_H(u)=[I:A G^u].
\]
Then $b_H$ is positive and nondecreasing, and
\[
 \varphi_H(\psi_G(u))=\int_0^u b_H(v)\,dv\qquad(u\geq0).
\]
Both sides of the change of parameter extend by the identity on
$[-1,0]$. No normality of $H$ is required.
\end{lemma}
\begin{proof}
For $u\geq0$, put $t=\psi_G(u)$ and $U=G_t=G^u$. Subgroup
compatibility identifies $H_t=H\cap U$, and $U\subset I$ is normal.
Consequently $AU$ is a subgroup of $I$. If $r=[A:A\cap U]$ and
$a=[I:AU]$, the subgroup index formulas give
\[
 r=[AU:U],\qquad ra=[I:U],\qquad
 r|H_t|=|A|,\qquad [I:U]|U|=|I|.
\]
Multiplying these equalities, and using that all the groups contain the
identity, proves
\[
 a=\frac{|H_t|}{|A|}\frac{|I|}{|U|}.
\]
This is exactly the product of the slopes of $\varphi_H$ at $t$ and
$\psi_G$ at $u$. Their finitely many exceptional arguments give only
finitely many exceptional $u$: besides the breaks of $\psi_G$, exclude
the images under $\varphi_G$ of the breaks of $\varphi_H$. The chain
rule therefore gives derivative $b_H(u)$ off a finite set. The composite
is continuous and vanishes at zero. Since $G^u$ decreases, so does
$AG^u$, and the finite index $b_H(u)$ is positive and nondecreasing.
It is thus integrable on every compact interval. The fundamental theorem
of calculus off a finite set gives the formula. For $-1\leq u\leq0$,
both finite Herbrand functions and their inverses are the identity,
which proves the stated extension.
\end{proof}
